% Modernised reading — fonds Alexandre Grothendieck, shelfmark 87
% (pages 1-40, the whole folder).
%
% This is an INTERPRETATION, not a transcription. It re-reads the pages in
% today's notation and vocabulary, states the hypotheses the manuscript leaves
% implicit, and drops the critical apparatus so that the mathematics reads
% straight through. Everything the brackets and underlines carried moves into a
% footnote. It is held to being mathematically correct as it stands; where the
% manuscript is loose or mistaken, the true statement is what appears here and
% the footnote says what the page has. In any dispute of substance the
% transcription governs: whoever cites Grothendieck must cite batch-01.fr.tex
% (pp. 1-20) or batch-02.fr.tex (pp. 21-40).
%
% Pass: Opus 5.5 (claude-opus-5-5), 2026-10-10 — first pass, unchecked against the pages by a human.
% The folder was read whole, its two batches together; this is one document
% for the shelfmark. It was made on Opus 5.5 and has not been measured against
% the reference reading of folder 115 (made on Opus 5). The literature named in
% the footnotes is cited from memory and was not checked against the sources.
%
% Scope. Page 1 is a blank wrapper, pages 16, 20 and 31 are blank; page 15 is
% a cover carrying « Cubes » (his) and « (13 p.) » (not obviously his). None
% of these carries mathematics.
%
% Spine.
%   pp. 2-14   I. Geometry of convex sets: programme (2); facettes = relative
%              interiors of faces, Fac(C), order, sups, extreme points,
%              Minkowski, finiteness (3-5); polars and Hahn-Banach (5-7);
%              « secteurs polyédraux » = polyhedra, Minkowski-Weyl (7-10);
%              polarity reverses the face order (10-12); « polyèdre
%              combinatoire strict », dimension function, diamond (12-14).
%   pp. 15-26  II. Cubes: intrinsic characterisation by the set of affine
%              spans of faces, axioms a)-e), product of 1-cubes, barycentric
%              model (17-19); « épures de n-cubes » (B with a free involution)
%              and the module they define, orientation torsor (21-22); faces of
%              the cube = partial sections (23-25); the real cube, W_n, the
%              equivalence convex cubes = épures (26).
%   pp. 27-30  III. Second mouture of the facettes (27-29), axiomatic
%              « polyèdre combinatoire » (lattice, chains, diamond) and the
%              census in dimensions -1 to 3 (29-30).
%   pp. 32-36  IV. Dimension 2 in rank 3: the « carte lisse combinatoire »,
%              flags and their involutions, oriented vertices/edges/faces;
%              axiom D unpacked; Sup/Inf lemma; degree >= 3 corollaries.
%   pp. 37-38  V. Finite posets whose lower intervals realise to spheres:
%              the dimension function delta, intervals.
%   pp. 39-40  VI. The diamond condition (Sigma), repères (flags), the
%              involutions sigma_i.
%
% Notation fixed for the whole document.
%  - « Facette » is kept, in his sense: the relative interior of a face (a
%    « face ouverte »), NOT a facet in today's sense (codimension one), which
%    is called « facette de codimension 1 » or « facette sous-maximale » here.
%    A « facette fermée » is a face. Fac(C) is the ordered set of facettes of C,
%    the empty one included (as on pp. 3-14; p. 27 excludes it, see there).
%  - 0 and 1 are the least and greatest elements of an ordered set.
%  - The polar is his: A° = {x' | <x,x'> >= -1 on A}, the opposite of today's
%    usual one. The correspondence of affine subspaces of p. 10, which he marks
%    with a small open ring distinct from °, is written F^vee here.
%  - « Polyèdre convexe » (his) = convex polytope; « secteur polyédral » (his)
%    = convex polyhedron, possibly unbounded. Both his words are kept.
%  - Cube: B a set of 2n elements with the free involution sigma (b |-> -b),
%    I = B/sigma; épure = (B, sigma).
%  - Rank-3 map (pp. 32-36): S, A, F vertices, edges, faces; vec A the
%    vertex-edge incidences, A^up the edge-face incidences; Dr the flags (he
%    writes A with a corner-and-arrow exponent); sigma_1, epsilon, sigma_2
%    change the vertex, the edge, the face of a flag.
%  - Fraktur X for the ordered sets of pp. 32-40, as in the transcription.
%
% Substantive departures from the pages, each footnoted where it occurs:
%  - p. 2: the three formulas for interior, boundary and closure are unions
%    over the lines through x, not intersections.
%  - p. 4, (5): « C' closed iff Phi is a down-set » holds when Fac(C) is finite
%    (polyhedral case), not in general (the disc).
%  - p. 6: the two corollary formulas are (Env(A1,A2))° = A1° n A2° and
%    (A1 n A2)° = Env(A1°, A2°); hypotheses supplied for the compactness
%    criteria.
%  - p. 11: the facettes paired by polarity are those with 0 not in E_F; the
%    page's « i.e. » writes 0 in E_F.
%  - p. 12: « polyèdre combinatoire strict » is undefined on the leaves; read
%    as « ordered set isomorphic to Fac(C), C a convex polytope ».
%  - p. 17: the injectivity in c) comes from a); b) gives the partial
%    sections.
%  - p. 26: B_n = {±e_i} is the set of extreme points of the POLAR of C_n
%    (the centres of the facets of C_n), not of C_n.
%  - p. 27: the definition of facette as written gives unions of facettes;
%    facettes are the minimal ones, i.e. the classes of R. The margin's
%    « finite intersection of half-spaces » needs compactness for (4)-(5).
%  - p. 28: closed faces of dim <= d-i are among the F n C, but not all such
%    F n C are faces (transversal line at the midpoint of an edge of a cube).
%  - p. 29: d) on p. 29-30 fails for atoms and coatoms; the diamond c) as
%    written (« there exists j' ») is too weak for the census of p. 30, which
%    needs exactly two middle elements, as Cor. 7 (p. 14) and (Sigma) (p. 39)
%    have.
%  - p. 32: the open intervals are spheres S^{delta-2}, not S^{delta-1}.
%  - p. 33: the labels of the arrows S~ -> S and F~ -> F seem exchanged.
%  - p. 34: the regrouping D = d1 + d2 + d'1 + d'2 loses one case (two faces
%    sharing two vertices and no edge); counterexample: a torus of 12 squares,
%    Z^2 / <(2,2),(3,-3)>. Checked by the edition with a short computation.
%  - p. 38: Lemma a) reads d - delta(x) - 2 for d - delta(x) - 1. More
%    seriously, the Proposition and Lemma are false for topological spheres
%    and manifolds in general (double suspension of a homology 3-sphere); what
%    holds is the homology version, and the sphere version for face lattices
%    of polytopes (by p. 12) and in low dimension.
%
% Modern names given and footnoted as ours: relative interior / face, face
% lattice, polar (dual) polytope, bipolar theorem, Hahn-Banach (his own),
% Minkowski (finite-dimensional Krein-Milman), Minkowski-Weyl, hyperoctahedral
% group, cross-polytope, diamond property, abstract polytope, thin poset,
% flag, combinatorial map, CW poset, homology sphere, double suspension.
%
% Announced and never established in the folder: the definition of
% « polyèdre combinatoire strict » (pp. 12, 14); the sequel to « Si J c I, on
% considère » (p. 19); the proof of the Prop. of p. 26 (« laissée comme
% exercice »); items (5)-(6) of p. 27, mostly illegible; the connectedness
% axioms wanted in dimension 3 (p. 30) and dimension 4 (heading only); the
% generalisation to « polyèdres lisses de dim quelconque » (p. 32); the proof
% of the Lemma of p. 38; the relations between sigma_i and sigma_{i+1} on the
% repères (pp. 39-40).
\documentclass[11pt,a4paper]{article}
\input{../preamble/grothendieck}

\folder{87}
\pages{1}{40}
\dating{s.d. — le groupe « Géométrie et topologie combinatoire » (69 à 102) est daté 1976-[vers 1986]}
\watermark{Édition de démonstration\\interprétation personnelle de l'œuvre}
\foldertitle{Polyèdres convexes : notes manuscrites (s.d.) — lecture modernisée du dossier entier}

\begin{document}

\section*{Résumé}

\begin{resume}
Un cube, une pyramide, un ballon de football à facettes : ce sont des
\emph{polyèdres convexes}, des solides sans creux, bordés de faces planes. Leur
bord se décompose en morceaux de toutes les dimensions --- des sommets, des
arêtes, des faces --- et l'on peut oublier la géométrie pour ne garder que la
façon dont ces morceaux s'emboîtent : tel sommet est au bout de telle arête,
telle arête borde telle face. Ces quarante pages demandent ce qui reste d'un
polyèdre quand on ne garde que cela, et si l'on peut reconnaître, dans un
tableau d'emboîtements écrit sur le papier, l'ombre d'un vrai polyèdre.

La première partie est un petit cours, très classique d'allure, sur les parties
convexes d'un espace réel. Grothendieck y découpe tout convexe en
« facettes » --- l'intérieur d'une face, sans son bord --- et range ces
facettes par l'ordre « l'une est dans l'adhérence de l'autre ». Il y introduit
l'outil décisif, la \emph{polarité} : à un polyèdre contenant l'origine en son
intérieur on associe un polyèdre dual, où les sommets deviennent des faces et
les faces des sommets, comme le cube et l'octaèdre s'échangent. La polarité
renverse l'ordre des facettes. Il en tire une propriété très simple et très
rigide, qu'on appelle aujourd'hui la propriété du \emph{losange} : entre une
facette et une autre de deux dimensions plus grande, il y a toujours exactement
deux facettes intermédiaires --- parce qu'un polyèdre de dimension 1 est un
segment, qui a deux bouts.

La deuxième partie s'arrête sur le cube. Un cube de dimension $n$ a $2n$ faces
de dimension $n-1$, rangées par paires de faces opposées, et Grothendieck
montre que c'est là toute sa structure : se donner un cube revient à se donner
un ensemble de $2n$ éléments muni d'une façon de les apparier --- ce qu'il
appelle une « épure de cube ». Les faces du cube se lisent alors comme des
choix partiels, dans certaines paires, d'un des deux éléments, et ses symétries
comme les permutations de l'ensemble qui respectent l'appariement.

La troisième partie retourne la question : quelles propriétés d'un ensemble
ordonné fini font qu'il « ressemble » à l'ensemble des facettes d'un polyèdre ?
Des bornes inférieures et supérieures, des chaînes de longueur constante, le
losange. Grothendieck essaie ces axiomes en petite dimension, voit qu'ils
laissent passer des objets qui ne sont pas des polyèdres --- deux triangles
disjoints, par exemple --- et étudie à fond le cas où le bord est une surface :
une \emph{carte}, avec ses sommets, ses arêtes, ses faces et ses drapeaux,
c'est-à-dire ses triplets sommet--arête--face emboîtés, sur lesquels opèrent
trois involutions. C'est la notion de carte des dossiers voisins 81 et 88,
et c'est dans le second que ces involutions engendrent le groupe
cartographique. Le dossier finit par
une version topologique --- un ensemble ordonné dont les intervalles, réalisés
géométriquement, sont des sphères --- et par une version purement
combinatoire, où la seule propriété du losange suffit à faire opérer des
involutions sur les drapeaux de toute longueur.

Ce qu'on voit à l'œuvre : un même résultat, la propriété du losange, est
obtenu trois fois, par la géométrie (page 14), par la topologie (page 38) et
posé comme axiome (page 39) ; et chaque fois qu'un axiome est écrit, il est
aussitôt mis à l'épreuve sur les plus petits exemples, jusqu'à ce que ceux-ci
disent ce qui manque. Ces notions ont aujourd'hui des noms --- treillis des
faces, polytopes abstraits, ensembles ordonnés CW --- que les pages ne
nomment pas.

\keywords{convex polytope, relative interior, face lattice, polar duality, bipolar theorem, Hahn-Banach theorem, Minkowski-Weyl theorem, diamond property, hypercube, hyperoctahedral group, cross-polytope, abstract polytope, graded lattice, combinatorial map, flag, CW poset, homology sphere}
\end{resume}

\pagerange{2}{40}
\section*{Le fil du dossier, et les conventions}

\subsection*{Les stations}

\begin{enumerate}
\item[I.] \emph{Géométrie des convexes} (pages 2 à 14), sous ce titre de sa
main. Le programme (page 2) ; les facettes d'un convexe, leur ordre, les
points extrémaux (pages 3 à 5) ; les polaires et le théorème de Hahn--Banach
(pages 5 à 7) ; les secteurs polyédraux (pages 7 à 10) ; polarité et facettes
(pages 10 à 12) ; la fonction dimension et le losange (pages 12 à 14).
\item[II.] \emph{Cubes} (pages 15 à 26), sous le titre de la couverture de la
page 15. Caractérisation intrinsèque d'un cube et reconstruction comme produit
de segments (pages 17 à 19) ; épures de $n$-cubes (pages 21 et 22) ; facettes
du cube (pages 23 à 25) ; le cube réel et son groupe (page 26).
\item[III.] \emph{Facettes, seconde rédaction, et polyèdres combinatoires}
(pages 27 à 30) : les facettes reprises sans la facette vide (pages 27 à 29),
puis une définition axiomatique du polyèdre combinatoire et le recensement des
petites dimensions (pages 29 et 30).
\item[IV.] \emph{La dimension 2} (pages 32 à 36) : le bord d'un polyèdre
combinatoire de dimension 3 comme « carte lisse combinatoire », ses drapeaux,
l'axiome D.
\item[V.] \emph{Ensembles ordonnés à intervalles sphériques} (pages 37 et 38).
\item[VI.] \emph{Le losange et les repères} (pages 39 et 40).
\end{enumerate}

La couverture de la page 15 porte aussi, au crayon et entourée, la mention
« (13 p.) », qui n'est pas sûrement de sa main ; c'est le nombre des feuillets
écrits 17 à 19 et 21 à 30, ce qui peut n'être qu'une coïncidence. Il n'y a
nulle part de pagination de sa main.

\subsection*{Deux rédactions des facettes, gardées distinctes}

Les facettes d'un convexe sont définies deux fois : page 3, comme classes d'une
relation d'équivalence, la facette vide étant admise et servant de plus petit
élément ; page 27, comme parties \emph{non vides}, avec une relation
d'équivalence formulée autrement. Les énoncés de la page 29 sur l'ensemble
ordonné des facettes, et les axiomes qui suivent, demandent de nouveau la
facette vide. On garde les deux rédactions séparées, et l'on dit à chaque fois
laquelle est en jeu ; $\mathrm{Fac}(C)$ contient toujours la facette vide.

\subsection*{Conventions}

\emph{Facettes.} Le mot est le sien, et on le garde, mais il faut le lire avec
soin : une \emph{facette} d'un convexe $C$ est ici l'intérieur relatif d'une
face de $C$ --- une « face ouverte » ---, de n'importe quelle dimension, et non
une face de codimension 1, que la langue d'aujourd'hui appelle
facette.\footnote{Le mot « facette » est celui de Bourbaki pour les
arrangements d'hyperplans (\emph{Groupes et algèbres de Lie}, chapitre V),
avec le même sens de partie ouverte dans son sous-espace affine ; le
rapprochement est le nôtre.} Une face de codimension 1 est dite ici
\emph{facette sous-maximale}, comme il le fait page 24. Une \emph{facette
fermée} est l'adhérence d'une facette dans $C$, c'est-à-dire une face.

\emph{Polyèdres.} Un \emph{polyèdre convexe}, chez lui, est compact : c'est ce
qu'on appelle aujourd'hui un polytope convexe. Un \emph{secteur polyédral} est
une intersection finie de demi-espaces fermés, bornée ou non : un polyèdre
convexe au sens d'aujourd'hui. On garde ses deux mots.

\emph{Ensembles ordonnés.} $0$ et $1$ désignent le plus petit et le plus grand
élément ; $y$ est un \emph{successeur} de $x$ si $x < y$ sans élément
intermédiaire (on dit aujourd'hui que $y$ couvre $x$) ; $\Phi_{x,y}$ est
l'intervalle fermé $\{z \mid x \leq z \leq y\}$ et $\mathfrak{X}_{x < \cdot < y}$
l'intervalle ouvert.

\emph{Polarité.} Le polaire est le sien,
$A^{\circ} = \{x' \mid \langle x, x' \rangle \geq -1 \ \forall x \in A\}$,
opposé de celui qu'on utilise d'ordinaire ($\leq 1$) ; les deux conventions
donnent les mêmes énoncés. La correspondance entre sous-espaces affines de la
page 10, qu'il marque d'un petit rond ouvert distinct du $^{\circ}$, est notée
ici $F \mapsto F^{\vee}$.

\emph{Cubes.} $B$ est un ensemble à $2n$ éléments muni d'une involution sans
point fixe $\sigma$, notée aussi $b \mapsto -b$, et $I = B/\sigma$ ; pour
$i \in I$, $B_i$ est la fibre, à deux éléments.

\emph{Cartes.} Aux pages 32 à 36, $S$, $A$, $F$ sont les sommets, arêtes et
faces ; $\vec{A} \subset S \times A$ et $A^{\uparrow} \subset A \times F$ les
incidences sommet--arête et arête--face ; $\mathrm{Dr} = \vec{A} \times_{A}
A^{\uparrow}$ l'ensemble des drapeaux\footnote{La page écrit $A$ avec un
exposant qui réunit la flèche verticale de $A^{\uparrow}$ et la flèche
horizontale de $\vec{A}$. La lettre $\mathfrak{X}$ des pages 32 à 40 rend un $X$
cursif.} ; $\sigma_{1}$, $\varepsilon$, $\sigma_{2}$ changent respectivement le
sommet, l'arête, la face d'un drapeau.

\subsection*{Ce que seule une lecture d'ensemble peut dire}

\begin{itemize}
\item Le « polyèdre combinatoire strict » des pages 12 et 14 n'est défini sur
aucun feuillet. La page 29 définit, sans le mot « strict », un polyèdre
combinatoire par des axiomes, et la page 30 montre aussitôt qu'en dimension 2
ces axiomes admettent des réunions disjointes de polygones, qui ne viennent
d'aucun polyèdre convexe. On lit donc « strict » comme « isomorphe à
l'ensemble ordonné des facettes d'un polyèdre convexe » --- la seule lecture
sous laquelle la démonstration de la page 12 a un sens --- et le polyèdre
combinatoire de la page 29 comme la notion plus large.
\item La propriété du losange --- exactement deux éléments entre deux éléments
de dimensions distantes de 2 --- est démontrée pour les polyèdres convexes
page 14 et page 29, déduite d'une hypothèse topologique page 38, et posée en
axiome page 39 sous le nom ($\Sigma$). L'axiome c) de la page 29 n'en est que
la moitié (« il existe »), et le recensement de la page 30 a besoin de l'autre.
\item Le modèle du cube par les coordonnées barycentriques, page 19, et le
module de la marge de la page 21 sont un même objet : le noyau de
$k^{B} \to k^{I}$.
\item Pour $n = 2$, l'épure de la page 21 est le « carré combinatoire » du
dossier 89 (page 13), lu sur les côtés plutôt que sur les sommets --- le carré
étant son propre dual. Les « contours » et « polygones » des pages 30 à 33 sont
les polygones combinatoires du dossier 89 (pages 19 et 20), connexité comprise
ou non ; les drapeaux et leurs involutions, pages 32, 33, 39 et 40, sont ceux du
dossier 88 (pages 6 à 8 et 10), où ils engendrent le groupe cartographique.
\end{itemize}

\pagerange{2}{14}
\section*{I. Géométrie des convexes (pages 2 à 14)}

\pagerange{2}{2}
\subsection*{Le programme (page 2)}

La page 2 est une table des matières : parties convexes d'un espace affine réel,
stabilité par intersection, enveloppe convexe décrite par les combinaisons
barycentriques ou par les simplexes, comparée au sous-espace affine engendré ;
stabilité par images directes et réciproques affines et par sommes ; sur une
droite, les convexes sont les intervalles, et une partie est convexe si et
seulement si sa trace sur toute droite est un intervalle ; dimension d'un
convexe. Puis l'intérieur. Soit $C$ un convexe de dimension finie qui
engendre $E$ affinement. Alors :

\begin{enumerate}
\item[a)] $\mathring{C} \neq \emptyset$ ;
\item[b)] si $x \in \mathring{C}$, et si $D$ parcourt les droites passant par
$x$,
\[
  \mathring{C} = \bigcup_{D} \operatorname{int}_{D}(D \cap C), \qquad
  \partial C = \bigcup_{D} \partial_{D}(D \cap C), \qquad
  \overline{C} = \bigcup_{D} \overline{D \cap C} ;
\]
\item[c)] si $C$ et $C'$ sont des convexes compacts qui engendrent $E$, et
$x \in \mathring{C}$, $x' \in \mathring{C}'$, il existe un unique
homéomorphisme $h : C \to C'$ tel que $h(x) = x'$, qui envoie chaque
demi-droite issue de $x$ sur la demi-droite parallèle et de même sens issue de
$x'$, et qui soit affine sur chacune.\footnote{Les trois formules de b) portent
sur la page le signe d'intersection ; ce sont des réunions --- une intersection
de parties de droites distinctes passant par $x$ se réduirait à $\{x\}$ ou à
rien. Pour la première, si $y$ est intérieur au segment $D \cap C$ de la droite
$D = (xy)$, $y$ est dans un segment ouvert $]x, z[$ avec $z \in C$, donc
intérieur à $C$ puisque $x$ l'est. Dans c), « parallèles » est une lecture
incertaine ; l'homéomorphisme est l'homéomorphisme radial défini par les
jauges de $C - x$ et $C' - x'$, et son unicité vient de ce qu'un homéomorphisme
envoie le bord sur le bord (invariance du domaine).}
\end{enumerate}

\pagerange{3}{5}
\subsection*{Facettes d'un convexe (pages 3 à 5)}

Soit $C$ un convexe d'un espace affine réel $E$. Pour $x, y \in C$, notons
$D_{x,y}$ la droite qui les joint et posons
\[
  R(x,y) \iff x = y \ \text{ou} \ x \text{ et } y \text{ sont intérieurs à }
  D_{x,y} \cap C \text{ dans } D_{x,y},
\]
c'est-à-dire : $x$ et $y$ appartiennent à un même segment ouvert contenu dans
$C$. C'est une relation d'équivalence. Une \emph{facette} de $C$ est une partie
$F$ telle que $F = \{y \in C \mid R(x,y)\}$ pour tout $x \in F$ : les facettes
sont les classes de $R$, et la partie vide, qui satisfait vacuement la
condition.\footnote{Une incise biffée et incertaine, « on exclut donc les $F \neq
\emptyset$ », suit la définition. La facette vide est utilisée aux pages 13 et
14 comme plus petit élément ($d(\emptyset) = -1$), et on la garde.} Les facettes
non vides sont exactement les intérieurs relatifs des faces non vides de
$C$ ; elles forment une partition de $C$.\footnote{Ce dictionnaire, et le mot
« face » au sens de Rockafellar (une partie convexe $G$ de $C$ telle que tout
segment de $C$ dont un point intérieur est dans $G$ y soit tout entier), sont
les nôtres.} On note $F_{x}$ la facette de $x$, $E_{F}$ le sous-espace affine
engendré par $F$, et $\mathrm{Fac}(C)$ l'ensemble des facettes.\footnote{Sur la
page, une abréviation lue « $\mathrm{Fac}(C)$ » dans tout le dossier ; l'indice
de $E_{C}$ est un $C$ appuyé.}

Un point $x$ est \emph{extrémal} si $\{x\}$ est une facette, c'est-à-dire si $x$
n'est intérieur à $D \cap C$ pour aucune droite $D$ passant par
$x$.

\begin{enumerate}
\item[(1)] Si $C$ est de dimension finie $d$, il a une unique facette de
dimension $d$, son intérieur relatif $\operatorname{int}_{E_{C}}(C)$.
\item[(2)] Pour $F \in \mathrm{Fac}(C)$, $E_{F} \cap C$ est une réunion de
facettes de $C$, qui sont les facettes du convexe $E_{F} \cap C$. En dimension
finie, $E_{F} \cap C = \overline{F} \cap C$ --- c'est la face de $C$ dont $F$
est l'intérieur relatif --- et $F = \operatorname{int}_{E_{F}}(E_{F} \cap C)$.
\item[(3)] Pour deux facettes $F$, $F'$, les conditions suivantes sont
équivalentes : $E_{F'} \subset E_{F}$ ; $E_{F'} \cap C \subset E_{F} \cap C$ ;
$F' \subset E_{F}$ ; et, en dimension finie, $F' \subset \overline{F}$ ;
$\overline{F'} \subset \overline{F}$. On note alors $F' \leq F$ : c'est une
relation d'ordre sur $\mathrm{Fac}(C)$, et pour $x, y \in C$,
\[
  F_{x} \leq F_{y} \iff x = y \ \text{ou} \ y \in
  \operatorname{int}_{D_{x,y}}(D_{x,y} \cap C).
\]
\end{enumerate}
En effet, $y$ intérieur au segment $D_{x,y} \cap C$ signifie $y \in \,]x,z[$
pour un $z \in C$, et alors $x$ est dans la face de $y$.\footnote{Un mot noirci
précède $F_{x}$ dans cette ligne ; la page écrit $\prec$.}

\begin{enumerate}
\item[(4)] Si $x = \sum \lambda_{i} x_{i}$ est une combinaison barycentrique
d'un nombre fini de points de $C$, à coefficients $\lambda_{i} > 0$, alors
$F_{x} = \sup_{i} F_{x_{i}}$ dans $\mathrm{Fac}(C)$. Les bornes supérieures
finies existent donc ; en fait toute famille a une borne supérieure, et
$\mathrm{Fac}(C)$ est un treillis complet, puisque les faces sont stables par
intersections quelconques.\footnote{La page conclut aux sup quelconques
« pourvu que $\dim C < +\infty$ » (deux mots incertains) ; l'hypothèse n'est
pas nécessaire. La stabilité des facettes fermées par intersections
quelconques est dite elle-même un peu plus bas sur la page.}
\item[(5)] Soit $C' \subset C$ une réunion de facettes, et
$\Phi = \{F \in \mathrm{Fac}(C) \mid F \subset C'\}$. Alors $C'$ est convexe si
et seulement si $\Phi$ est stable par sup finis. Si $C'$ est fermé dans $C$,
$\Phi$ contient, avec une facette, toutes celles qui sont plus petites ; la
réciproque est vraie quand $\mathrm{Fac}(C)$ est fini.\footnote{La page énonce
l'équivalence sans restriction (« car $\dim E < \infty$ », les mots « relat. »
et « majorées » étant incertains). Elle est fausse en général : dans un disque
fermé, la réunion des points d'un arc ouvert du cercle est une réunion de
facettes (des points extrémaux), $\Phi$ est stable par descente, et l'arc n'est
pas fermé. Elle est vraie si $\mathrm{Fac}(C)$ est fini, $C'$ étant alors une
réunion finie de faces, qui sont fermées dans $C$.}
\item[(6)] (En dimension finie.) $C'$ est de la forme $\overline{F} \cap C$,
$F \in \mathrm{Fac}(C)$, si et seulement si c'est une réunion de facettes,
convexe et fermée dans $C$ ; en termes de $\Phi$ : si et seulement si $\Phi$ a
un plus grand élément $F$ et est l'ensemble des facettes $\leq F$. Les facettes
fermées sont stables par intersections quelconques.\footnote{La page écrit
« que $\Phi$ soit fermé, convexe et réunion de facettes », mêlant $\Phi$ et
$C'$, et la numérotation hésite entre (5) et (6). La démonstration de
l'équivalence est la suivante : si $C'$ est une telle réunion et $x$ un point de
son intérieur relatif, tout $y \in C'$ est le bout d'un segment de $C'$ dont $x$
est un point intérieur, donc $C' \subset \overline{F_{x}} \cap C$ ; l'inclusion
inverse vient de ce que $C'$ est fermé dans $C$ et contient $F_{x}$. La même
caractérisation revient page 29, c).}
\item[(7)] Si $C$ est compact de dimension finie, $C$ est l'enveloppe convexe de
l'ensemble de ses points extrémaux.\footnote{C'est le théorème de Minkowski,
cas de la dimension finie du théorème de Krein--Milman, où l'adhérence est
inutile ; les noms ne sont pas sur la page.}
\item[(8)] Si $C$ est une partie convexe polyédrale de $E$ de dimension finie
--- réunion finie de simplexes, ou, ce qui revient au même, d'enveloppes
convexes de parties finies --- alors $\mathrm{Fac}(C)$ est fini.
\end{enumerate}

\textbf{Corollaire.} Les \emph{polyèdres convexes} sont les convexes compacts
$C$ tels que $\mathrm{Fac}(C)$ soit fini, ou, ce qui revient au même, tels que
l'ensemble $\mathrm{Fac}_{0}(C)$ de leurs points extrémaux soit fini.

\pagerange{5}{7}
\subsection*{Polaires et théorème de Hahn--Banach (pages 5 à 7)}

Soient $E$ et $E'$ deux espaces vectoriels réels de dimension finie en dualité
par $\langle\, ,\rangle$.\footnote{La page dit seulement « espaces vectoriels
duals l'un de l'autre » ; toute la suite est en dimension finie, et c'est
l'hypothèse sous laquelle le théorème de Hahn--Banach est démontré page 7. En
dimension infinie, il faudrait des topologies compatibles avec la dualité.}
Pour $A \subset E$, posons
\[
  A^{\circ} = \{x' \in E' \mid \langle x, x' \rangle \geq -1 \ \ \forall x \in A\}
  = \bigcap_{x \in A} H(x), \qquad
  H(x) = \{x' \mid \langle x, x' \rangle \geq -1\},
\]
et de même pour les parties de $E'$.

\textbf{Théorème 1 (des bipolaires).} Pour $A \subset E$, $A^{\circ\circ}$ est
l'enveloppe convexe fermée de $A \cup \{0\}$.

\textbf{Corollaires.} $A \mapsto A^{\circ}$ est une bijection, qui renverse
l'ordre, de l'ensemble des convexes fermés de $E$ contenant $0$ sur l'ensemble
analogue de $E'$, d'inverse $A' \mapsto A'^{\circ}$. Pour de tels convexes,
\[
  \bigl(\overline{\mathrm{Env}}(A_{1} \cup A_{2})\bigr)^{\circ} = A_{1}^{\circ}
  \cap A_{2}^{\circ}, \qquad
  (A_{1} \cap A_{2})^{\circ} = \overline{\mathrm{Env}}(A_{1}^{\circ} \cup
  A_{2}^{\circ}),
\]
où $\overline{\mathrm{Env}}$ désigne l'enveloppe convexe fermée ; et
\[
  0 \in \operatorname{int}(A) \iff A^{\circ} \text{ est compact}, \qquad
  A \text{ est compact} \iff 0 \in \operatorname{int}(A^{\circ}).
\]
Ces deux critères supposent $A$ convexe fermé contenant $0$.\footnote{Sur la page, aucun signe de polaire sur le membre de gauche de la
première formule ni sur les arguments du second membre de la seconde. Les
deux critères de compacité sont énoncés sans hypothèse ; l'hypothèse est
fournie par nous.}

Le théorème 1 équivaut au suivant :\footnote{La ligne, marquée d'un grand
crochet, est en partie biffée et un mot en est illisible : « Le théorème est
conséquence [\ldots] est équivalent ».}

\textbf{Théorème 1'.} Une partie de $E$ est convexe et fermée si et seulement si
elle est intersection de demi-espaces fermés.

Il résulte du

\textbf{Théorème 1'' (Hahn--Banach).} Soient $U$ une partie convexe ouverte de
$E$ et $F \subsetneq E$ un sous-espace vectoriel tel que $F \cap U =
\emptyset$. Il existe un hyperplan $H \supset F$ tel que $H \cap U =
\emptyset$.

\emph{Démonstration (page 7).} Remplaçant $U$ par le cône $\bigcup_{\lambda > 0}
\lambda U$, qui est encore convexe, ouvert et disjoint de $F$, on peut supposer
que $U$ est un cône ouvert. Il suffit, si $F$ n'est pas un hyperplan, de trouver
$F' \supset F$ avec $\dim F' = \dim F + 1$ et $F' \cap U = \emptyset$, puis
d'itérer. Pour cela on se place dans un sous-espace $G \supset F$ de dimension
$\dim F + 2$, et, passant au quotient par $F$, on est ramené au cas où
$\dim E = 2$ et $F = 0$, avec $U \neq \emptyset$. Un cône convexe ouvert du
plan qui ne contient pas l'origine est alors ou bien un demi-plan ouvert, ou
bien un secteur $\{\lambda_{1} e_{1} + \lambda_{2} e_{2} \mid \lambda_{1},
\lambda_{2} > 0\}$ pour une base $(e_{1}, e_{2})$ ; dans les deux cas une droite
de son bord convient.\footnote{La page saute de « il suffit de prouver que
$\exists F' \supset F$ […] $\dim F' = \dim F + 1$ » à la suite sans écrire
$F' \cap U = \emptyset$, qu'on rétablit, et ne dit pas comment on trouve la
droite (« Dans les deux cas on gagne… »). Le même mot, illisible, ouvre deux
phrases de la démonstration.}

\pagerange{7}{10}
\subsection*{Secteurs polyédraux (pages 7 à 10)}

\textbf{Proposition.} Soit $C$ une partie d'un espace affine $E$ de dimension
finie. Conditions équivalentes :
\begin{enumerate}
\item[a)] $C$ est une intersection finie de demi-espaces fermés ;
\item[b)] $C$ est l'enveloppe convexe fermée d'un ensemble fini de points et de
demi-droites fermées, qu'on peut prendre toutes issues d'un point $x \in C$
donné à l'avance ;
\item[b')] (si $C$ est compact) $C$ est un polyèdre convexe.
\end{enumerate}

C'est le théorème de Minkowski--Weyl.\footnote{Le nom est le nôtre. La
précision « issues d'une origine » est d'une lecture incertaine. L'adhérence
est nécessaire dans b) quand les demi-droites ont une même origine :
l'enveloppe convexe de $\{v\}$ et d'une demi-droite issue de $x \neq v$ ne
contient pas la demi-droite parallèle issue de $v$, que contient son
adhérence.}

\emph{Démonstration (pages 8 et 9).} Si $C = \emptyset$, il n'y a rien à
démontrer. Les deux conditions ne changent pas quand on remplace $E$ par le
sous-espace affine $E_{C}$ engendré par $C$ ; on peut donc supposer que $C$ est
convexe fermé et engendre $E$, et prendre pour origine un point de
$\mathring{C}$. On pose $C' = C^{\circ}$ ; alors $C = C'^{\circ}$, et $C'$ est
compact.\footnote{Un premier essai, encadré et barré, où se lisent « a)
$\Rightarrow$ b) », « b) $\Rightarrow$ a) » et « on gagne par polarité », est
remplacé par ce qui suit. « $C$ engendre $E$ » est une lecture incertaine.}

\emph{b) $\Rightarrow$ a).} Le polaire d'un point $x$ est le demi-espace $H(x)$,
celui d'une demi-droite $\mathbf{R}_{+} d$ issue de l'origine le demi-espace
$\{x' \mid \langle d, x'\rangle \geq 0\}$, et le polaire d'une réunion est
l'intersection des polaires. Donc $C'$ est une intersection finie de
demi-espaces fermés ; étant compact, c'est un polyèdre convexe par le lemme
ci-dessous, c'est-à-dire l'enveloppe convexe d'un ensemble fini de $E'$, et
$C = C'^{\circ}$ est l'intersection finie des polaires de ses points : a).

\emph{a) $\Rightarrow$ b).} Par polarité, $C'$ est l'enveloppe convexe fermée d'un
ensemble fini de points ; par ce qui précède, appliqué dans $E'$ à $C'$ avec une
origine prise dans son intérieur relatif, c'est une intersection finie de
demi-espaces fermés, qu'on peut écrire, puisque $0 \in C'$, sous les formes
$\langle y_{j}, x'\rangle \geq -1$ et $\langle d_{k}, x'\rangle \geq 0$. Donc
$C = C'^{\circ}$ est l'enveloppe convexe fermée des points $y_{j}$ et des
demi-droites $\mathbf{R}_{+} d_{k}$ issues de l'origine : b), sous la forme
précisée.\footnote{La page dit seulement « d'après ce qui précède » ; le
recentrage de l'origine, nécessaire pour appliquer b) $\Rightarrow$ a) à $C'$,
est explicité par nous.}

\textbf{Lemme (page 9).} Une intersection finie de demi-espaces fermés n'a
qu'un nombre fini de points extrémaux ; c'est donc un polyèdre convexe si elle
est compacte.\footnote{La page met « compact » dans l'hypothèse (lecture
incertaine) ; l'énoncé et la démonstration n'en ont pas besoin.}

Par récurrence sur $\dim E$, puis sur le nombre $\nu$ de demi-espaces, le cas
$\nu \leq 1$ étant trivial. On est ramené à $C = C_{1} \cap H$, où $C_{1}$ n'a
qu'un nombre fini de points extrémaux et $H$ est un demi-espace fermé limité par
l'hyperplan $H_{0}$. Pour $x \in C$ : si $x \notin H_{0}$, $x$ est extrémal dans
$C$ si et seulement s'il l'est dans $C_{1}$, la propriété étant locale ; si
$x \in H_{0}$, $x$ est extrémal dans $C$ si et seulement s'il l'est dans
$H_{0} \cap C$, car un segment de $C$ qui a un point intérieur sur $H_{0}$ est
contenu dans $H_{0}$. Or $H_{0} \cap C$ est une intersection finie de
demi-espaces de $H_{0}$, de dimension plus petite.\footnote{La page note $\nu$
d'après une ligne biffée, « sur $\nu = n + \dim E$ » ; un croquis en marge
montre une bande entre deux droites parallèles.}

\textbf{Définition.} Un \emph{secteur polyédral} est une partie satisfaisant aux
conditions équivalentes de la proposition --- un polyèdre convexe au sens
d'aujourd'hui, borné ou non.

\textbf{Proposition (page 10).} Si $E$ et $E'$ sont de dimension finie et en
dualité, $A \mapsto A^{\circ}$ et $A' \mapsto A'^{\circ}$ sont des bijections
inverses l'une de l'autre entre secteurs polyédraux de $E$ contenant l'origine
et secteurs polyédraux de $E'$ contenant l'origine.

\pagerange{10}{12}
\subsection*{Polarité et facettes (pages 10 à 12)}

Soient $E$, $E'$ de dimension $n$ en dualité. À un sous-espace affine non vide
$F$ de $E$ ne contenant pas l'origine on associe
\[
  F^{\vee} = \{x' \in E' \mid \langle x, x'\rangle = -1 \ \ \forall x \in F\},
\]
sous-espace affine de $E'$ ne contenant pas l'origine, et l'on a
\[
  \dim F + \dim F^{\vee} = n - 1.
\]
C'est une bijection involutive, caractérisée par deux propriétés : elle renverse
l'inclusion, et à un point $x \neq 0$ elle associe l'hyperplan
$\{x\}^{\vee} = \partial H(x)$, bord du demi-espace polaire de
$x$.\footnote{L'exposant de $F^{\vee}$ est sur la page un petit rond ouvert,
tracé autrement que le $^{\circ}$ des polaires, et le bord est noté par un point
en exposant, comme $\dot{C}$ page 2.}

Une \emph{variété d'appui} d'un convexe $C$ est un sous-espace affine de la forme
$E_{F}$, $F \in \mathrm{Fac}(C)$ ; les variétés d'appui sont en bijection avec
les facettes, et cette bijection respecte l'ordre.

\textbf{Théorème.} Soient $C \subset E$ et $C' \subset E'$ deux secteurs
polyédraux polaires l'un de l'autre. Il existe une unique application
$\varphi$ de l'ensemble des facettes non vides $F$ de $C$ telles que
$0 \notin \overline{F}$ --- c'est-à-dire, puisque $0 \in C$, telles que
$0 \notin E_{F}$ --- dans l'ensemble analogue pour $C'$, telle que
\[
  E_{\varphi(F)} = (E_{F})^{\vee}.
\]
Elle est bijective et renverse l'ordre.\footnote{Les deux « i.e. » de la page
écrivent $0 \in E_{F}$ et $0' \in E_{F'}$ ; c'est $\notin$ qu'il faut, pour que
$(E_{F})^{\vee}$ soit défini. L'équivalence entre $0 \notin \overline{F}$ et
$0 \notin E_{F}$ vient de $E_{F} \cap C = \overline{F} \cap C$ (page 3, (2)).
Exemple : pour $C = [-1, +\infty[$ dans $\mathbf{R}$, $C' = [0, 1]$, et
$\varphi$ envoie $\{-1\}$ sur $\{1\}$ ; la facette $\{0\}$ de $C'$, dont la
variété d'appui contient l'origine, traduit le fait que $C$ n'est pas borné.}

\textbf{Corollaire 1.} Un secteur polyédral n'a qu'un nombre fini de facettes :
prendre l'origine dans son intérieur relatif, et appliquer la dualité, le
polaire étant alors compact.

\textbf{Corollaire 2.} Si $C$ et $C'$ sont compacts, c'est-à-dire si l'origine
est intérieure à l'un et à l'autre, $\varphi$ se prolonge, en échangeant la
facette vide et la facette maximale, en un anti-isomorphisme d'ensembles
ordonnés $\mathrm{Fac}(C) \simeq \mathrm{Fac}(C')$.

C'est la dualité des polytopes : le polaire d'un polytope contenant l'origine
en son intérieur est un polytope, et son treillis des faces est l'opposé de
celui du premier.\footnote{Les noms « polytope dual », « treillis des faces »
sont les nôtres.}

\pagerange{12}{14}
\subsection*{Polyèdres combinatoires stricts et fonction dimension (pages 12 à 14)}

Appelons \emph{polyèdre combinatoire strict} un ensemble ordonné isomorphe à
$\mathrm{Fac}(C)$ pour un polyèdre convexe $C$, et notons $\Phi^{\circ}$
l'ensemble ordonné opposé de $\Phi$.\footnote{La notion n'est définie sur aucun
feuillet du dossier ; la définition est notre lecture. Elle est la seule sous
laquelle « a) résulte du Cor.~2 » a un sens, et la page 29 donne, sans le mot
« strict », une définition axiomatique qui laisse passer des objets qui ne sont
pas de cette forme (page 30).}

\textbf{Corollaire 3.} Soit $\Phi$ un polyèdre combinatoire strict.
\begin{enumerate}
\item[a)] $\Phi^{\circ}$ est un polyèdre combinatoire strict.
\item[b)] Pour $x \leq y$ dans $\Phi$, l'intervalle $\Phi_{x,y}$ est un polyèdre
combinatoire strict ; si $\Phi = \mathrm{Fac}(C)$, c'est $\mathrm{Fac}(K)$ pour
un polyèdre convexe $K$ de dimension $\dim y - \dim x - 1$.
\end{enumerate}

a) résulte du corollaire 2. Pour b) : $\Phi_{0,x} = \mathrm{Fac}(\overline{x})$
est strict ; donc aussi, par dualité, $\Phi_{x,1} = \bigl((\Phi^{\circ})_{0,x}
\bigr)^{\circ}$ ; et $\Phi_{x,y}$ est l'intervalle $[x, 1]$ de $\Phi_{0,y}$. La
dimension se lit sur la formule $\dim F + \dim F^{\vee} = n - 1$.\footnote{La
dimension de $K$ n'est pas sur la page 12 ; elle est donnée page 29, 2°, où
$\delta = \dim F'' - \dim F' - 1$. C'est le résultat classique selon lequel les
intervalles d'un treillis des faces de polytope sont des treillis des faces de
polytopes (figures de faces).}

Suivent deux essais, encadrés et barrés, d'un corollaire sur l'existence de
facettes de toutes dimensions, dont l'énoncé revient comme corollaire 6.

\textbf{Corollaire 4.} Soient $C$ un polyèdre convexe et $F_{1} < F_{2}$ dans
$\mathrm{Fac}(C)$. Pour que $F_{2}$ soit un successeur de $F_{1}$, il faut et il
suffit que $\dim F_{2} = \dim F_{1} + 1$.

La dimension étant strictement croissante sur $\mathrm{Fac}(C)$, la condition
est suffisante. Pour la nécessité, on se ramène par le corollaire 3 à
$F_{1} = 0$ (la facette vide) et $F_{2} = 1$ (la facette maximale), et
l'assertion devient : un polyèdre convexe non vide sans autre facette que la
vide et la maximale est un point. En effet, s'il est de dimension $\geq 1$, ses
points extrémaux, qui existent par le théorème de Minkowski, sont des facettes
intermédiaires.\footnote{Plusieurs mots de cet argument sont incertains :
« dualité », le signe $=$ dans « $\dim C = 0$ ou $-1$ », « vide »,
« trivial ». La page passe aussi par la facette maximale avant de passer aux
polaires.}

\textbf{Corollaire 5.} La fonction $d : F \mapsto \dim F$ sur $\mathrm{Fac}(C)$
est l'unique fonction telle que $d(F_{2}) = d(F_{1}) + 1$ quand $F_{2}$ est un
successeur de $F_{1}$, et $d(\emptyset) = -1$. Elle est strictement croissante.

\textbf{Corollaire 6.} Pour $F \in \mathrm{Fac}(C)$ et $-1 \leq i \leq \dim F$,
il existe $F' \leq F$ de dimension $i$.

C'est formel à partir du corollaire 5 : une chaîne de successeurs
$\emptyset = x_{0} < x_{1} < \cdots < x_{m} = F$, qui existe dans un ensemble
ordonné fini, vérifie $d(x_{i}) = i - 1$.\footnote{Une démonstration par
récurrence sur $\dim F$ est encadrée et barrée ; « Formel » est une lecture
incertaine.}

\textbf{Corollaire 7 (le losange).} Soient $\Phi$ un polyèdre combinatoire
strict, $d$ sa fonction dimension, et $x < y$ avec $d(y) = d(x) + 2$. Alors il y
a exactement deux éléments $z$ tels que $x < z < y$.

En effet, $\Phi_{x,y}$ est l'ensemble des facettes d'un polyèdre convexe de
dimension 1, c'est-à-dire d'un segment, qui a exactement deux extrémités. C'est
ce qu'on appelle aujourd'hui la propriété du losange (\emph{diamond
property}).\footnote{Le nom est le nôtre ; il est d'usage dans la théorie des
polytopes abstraits (Danzer, McMullen et Schulte), qui prend précisément pour
axiomes des ensembles ordonnés gradués ayant cette propriété.}

\pagerange{17}{19}
\section*{II. Cubes (pages 15 à 26)}

\subsection*{Caractérisation intrinsèque d'un cube (pages 15 à 19)}

La couverture de la page 15 porte le seul mot « Cubes ». Les pages 17 à 19
cherchent à caractériser un $n$-cube d'un espace affine $E$ sur un corps $k$ par
l'ensemble $\Phi$ des sous-espaces affines engendrés par ses
faces.\footnote{« $n$-cube » et la notation $\mathrm{Espaff}(E)$ de l'ensemble
des sous-espaces affines sont des lectures incertaines.} On note $\Phi^{**}$
l'ensemble $\Phi$ privé de son plus petit et de son plus grand
élément,\footnote{L'exposant est fait de deux petits signes, lus $**$.} et $B$
l'ensemble des éléments maximaux de $\Phi^{**}$. Les axiomes :

\begin{enumerate}
\item[a)] $\Phi$ contient $\emptyset$ et $E$, qui sont donc son plus petit et son
plus grand élément, et tout élément de $\Phi^{**}$ est intersection d'éléments
de $B$.
\item[b)] Pour tout $b \in B$ il existe un unique $b' \in B$ tel que $b$ et $b'$
n'aient pas de minorant commun dans $\Phi^{**}$, c'est-à-dire que $b \cap b'$
ne contienne aucun élément non vide de $\Phi$. On obtient ainsi une involution
$\sigma : b \mapsto b'$ de $B$, sans point fixe ; soit $I = B/\sigma$.
\item[c)] Pour $V \in \Phi \setminus \{\emptyset\}$, l'ensemble $B_{V}$ des
$b \in B$ qui majorent $V$ ne contient jamais à la fois $b$ et $\sigma b$ (par
b)) : c'est une section partielle de $B \to I$. L'application
\[
  \Phi \setminus \{\emptyset\} \longrightarrow \Gamma_{\mathrm{part}}(B/I),
  \qquad V \longmapsto B_{V},
\]
qui est injective par a), est bijective.
\item[d)] $n = \frac{1}{2}\operatorname{card} B$ est la dimension combinatoire de
$\Phi$ et la dimension de $E$, et pour tout $V \in \Phi$ la dimension
combinatoire de $V$ dans $\Phi$ est sa dimension affine.
\item[e)] Si $b \in B$, alors $b \cap \sigma b = \emptyset$ : par d), $b$ et
$\sigma b$ sont deux hyperplans parallèles.
\end{enumerate}

Une note en marge dit que a), b), c) expriment que l'ensemble ordonné $\Phi$ est
un \emph{cube combinatoire} : l'ensemble des sections partielles de $B$ sur $I$,
ordonné par l'inverse du prolongement, augmenté d'un plus petit
élément.\footnote{La page écrit « injective par b) » ; l'injectivité vient de
a), puisque $V$ est l'intersection des éléments de $B_{V}$, et b) donne que
$B_{V}$ est une section partielle. Le dernier tiers de la page 17, serré et en
surcharge, n'est lu que par fragments.} Une remarque entre crochets explique
pourquoi e) est nécessaire : a) implique que la borne inférieure dans $\Phi$
d'une partie minorée de $\Phi \setminus \{\emptyset\}$ est son intersection,
mais non qu'une partie dont la borne inférieure est $\emptyset$ ait une
intersection vide. Exemple : un quadrilatère qui n'est pas un parallélogramme,
dont deux côtés opposés se coupent hors de lui.\footnote{« Quadrilatère » est
une lecture incertaine, appuyée sur le croquis.}

\emph{Reconstruction.} Pour $i \in I$, les deux hyperplans de la fibre de
$B \to I$ en $i$ ont une même direction $\mathbb{H}_{i}$, hyperplan de l'espace
des translations $T$ de $E$. Ces $n$ hyperplans sont indépendants --- un sommet
du cube, élément minimal non vide de $\Phi$, est de dimension $0$ par d) et
s'écrit comme intersection de $n$ hyperplans, un par fibre ---, d'où une
décomposition de $T$ en produit de droites
\[
  T \simeq \prod_{i \in I} T_{i}, \qquad T_{i} = T/\mathbb{H}_{i},
\]
et une décomposition du $T$-torseur $E$ en produit de $T_{i}$-torseurs
$E_{i} = E/\mathbb{H}_{i}$, droites affines. Chaque $E_{i}$ contient l'image
$B_{i}$ de la fibre en $i$, deux points distincts.\footnote{La page écrit
« l'espace des translations $T$ de $V$ » ; c'est $E$. La lettre rendue
$\mathbb{H}$ est un $H$ barré de traits verticaux. L'indépendance des
directions n'est pas justifiée sur la page.} Or, $i$ étant fixé, la donnée d'un
triplet $(T_{i}, E_{i}, B_{i} \subset E_{i})$ --- une droite affine et deux
points distincts --- équivaut à celle de l'ensemble à deux éléments $B_{i}$, vu
comme torseur sous $\{\pm 1\}$ : on prend
\[
  T_{i} = \operatorname{Ker}\bigl(k^{B_{i}} \xrightarrow{\ \varepsilon\ } k\bigr),
  \qquad E_{i} = \varepsilon^{-1}(1),
\]
où $\varepsilon$ est la somme des coordonnées, $B_{i}$ s'envoyant sur les deux
fonctions de Dirac. Ce sont les coordonnées barycentriques sur une droite
relativement à deux points. Le $n$-cube est ainsi le produit de $n$ cubes de
dimension 1, et l'on a globalement
\[
  T = \operatorname{Ker}\bigl(k^{B} \xrightarrow{\ \mathrm{tr}\ } k^{I}\bigr),
  \qquad E = \mathrm{tr}^{-1}(1),
\]
où $\mathrm{tr}$ somme les coordonnées sur chaque fibre et $1$ est la fonction
constante. La donnée d'une famille $(T_{i}, E_{i}, B_{i})_{i \in I}$ équivaut
donc à celle de l'ensemble $B$ muni de $\sigma$.\footnote{La page dit qu'elle
équivaut « à la donnée d'une [mot illisible] de $n$-cube » ; c'est l'objet que
la page 21 appellera épure de $n$-cube, mais on ne prétend pas lire le mot.}
La page s'interrompt sur « Si $J \subset I$, on considère ».

\pagerange{21}{22}
\subsection*{Épures de $n$-cubes (pages 21 et 22)}

Soit $k$ un anneau commutatif où $2$ est inversible.\footnote{« comm. » est
récrit au-dessus de la ligne ; la lecture est incertaine.} On considère deux
groupoïdes.
\begin{itemize}
\item $\mathcal{C}_{n}$ : les $k$-modules $E$ munis d'une décomposition en somme
directe $E = \bigoplus_{i \in I} E_{i}$, $\operatorname{card} I = n$, et, pour
chaque $i$, d'une partie $B_{i} \subset E_{i}$ à deux éléments, telle que chaque
$b \in B_{i}$ soit une base de $E_{i}$ et que $B_{i} = -B_{i}$.
\item $\mathcal{D}_{n}$ : les ensembles $B$ à $2n$ éléments munis d'une
involution sans point fixe --- d'une action libre de $\{\pm 1\}$ ---, que la
marge appelle \emph{épures de $n$-cubes}.
\end{itemize}
On passe de $\mathcal{C}_{n}$ à $\mathcal{D}_{n}$ par $B = \bigcup_{i} B_{i}$,
$\sigma(b) = -b$, et en sens inverse par
\[
  E = \bigoplus_{i \in I} B_{i} \wedge_{\{\pm 1\}} k ,
\]
où $\{\pm 1\}$ opère sur $k$ par homothéties et $\wedge$ désigne le produit
contracté : $B_{i} \wedge_{\{\pm 1\}} k$ est un module libre de rang 1 dont $B_{i}$
est l'ensemble de deux générateurs opposés. Les deux groupoïdes sont donc
équivalents.\footnote{La condition $\operatorname{card} B_{i} = 2$ demande
seulement $b \neq -b$, c'est-à-dire $2 \neq 0$ dans $k$ ; l'inversibilité de $2$
sert au paragraphe suivant. La page note les deux catégories $C_{n}$ et
$\mathcal{D}_{n}$ ; on écrit $\mathcal{C}_{n}$ pour éviter la confusion avec le
cube $C_{n}$ de la page 26. Le mot « cube » de la marge est incertain, et
« combinatoire » y est biffé.}

La marge donne une autre description : $E$ est le sous-module de $k^{(B)}$ formé
des $\sum_{b} \lambda_{b} e_{b}$ avec $\lambda_{-b} = -\lambda_{b}$, d'où la suite
exacte
\[
  0 \to E \to k^{B} \xrightarrow{\ \mathrm{tr}\ } k^{I} \to 0 ,
\]
$b \in B$ correspondant à $e_{b} - e_{-b}$. C'est exactement l'espace $T$ des
translations du cube de la page 19 ; et quand $2$ est inversible, la bijection
$\lambda \mapsto \mu$, $\mu_{b} = \lambda_{b} - \lambda_{-b}$, d'inverse
$\lambda_{b} = \frac{1}{2}(1 + \mu_{b})$, identifie le modèle barycentrique
$\mathrm{tr}^{-1}(1)$ de la page 19 au modèle antisymétrique des pages 23 et
26.\footnote{Ce rapprochement des deux modèles est le nôtre ; il explique
l'hypothèse « $2$ inversible ».}

\emph{Structure euclidienne et orientation (page 22).} $E$ porte une forme
bilinéaire canonique, somme orthogonale des formes $(x, y) \mapsto xy$ sur les
$B_{i} \wedge k$, pour laquelle chaque $b \in B$ est de carré 1.\footnote{« de
produit scalaire » est un ajout interlinéaire incertain, et « bilinéaire » une
lecture incertaine.} Sur $k = \mathbf{R}$, notons $\omega_{E}$ l'ensemble à deux
éléments des orientations de $E$ et $\omega_{I}$ celui des orientations de
l'ensemble fini $I$ (ses ordres totaux à permutation paire près). On a un
isomorphisme canonique de $\{\pm 1\}$-torseurs
\[
  \omega_{E} \simeq \omega_{I} \wedge \bigwedge_{i \in I} B_{i} .
\]
En effet, la formule d'associativité donne $\omega_{E} \simeq \omega_{I} \wedge
\bigwedge_{i} \omega_{E_{i}}$, et $B_{i}$, ensemble de deux bases opposées de la
droite $E_{i}$, s'envoie injectivement, donc bijectivement, dans
$\omega_{E_{i}}$.\footnote{Le mot « sur $\mathbf{R}$ » est un petit « $R$ » écrit
sous la ligne après « bases » ; c'est lui qui fixe le cadre réel de ce
paragraphe. La marge porte « [\ldots] vérifier : ». En bas de page, souligné comme
un titre : « épure de $n$-cube orienté ».} Une \emph{épure orientée} est une
épure munie d'un élément de ce torseur.

\pagerange{23}{25}
\subsection*{Facettes du cube (pages 23 à 25)}

Le \emph{polyèdre combinatoire} associé à un polyèdre convexe $C$ de dimension
$n$ est l'ensemble de ses facettes, ordonné par $F \leq F' \iff F \subset
\overline{F'}$.\footnote{Un ajout interlinéaire, « $\neq$ l'unique facette
maxim. ? », semble exclure la facette maximale, ce que fait la liste
$\Phi_{0}, \ldots, \Phi_{n-1}$ qui suit.} La dimension d'une facette s'y lit
comme dimension combinatoire : la longueur commune des chaînes maximales
$F_{0} < F_{1} < \cdots < F_{m} = F$ partant d'un sommet. Si $\Phi_{i}$ est
l'ensemble des facettes de dimension $i$, l'ordre se reconstitue à partir des
seules relations d'incidence entre $\Phi_{i}$ et $\Phi_{i+1}$, puisque toute
relation $F < F'$ se raffine en une chaîne de successeurs.\footnote{La page
écrit $0 \leq i \leq n$ pour les relations d'incidence, alors que la liste
s'arrête à $\Phi_{n-1}$.}

Le cube réel associé à une épure est
\[
  C(B, \sigma) = \Bigl\{ \lambda = \sum_{b \in B} \lambda_{b} e_{b} \Bigm|
  \lambda_{-b} = -\lambda_{b},\ |\lambda_{b}| \leq 1 \Bigr\}
\]
dans le module $E$ ci-dessus, pour $k = \mathbf{R}$.

\textbf{Proposition.} Les facettes non vides du cube $C(B, \sigma)$ sont en
bijection avec les sections partielles $s : J \to B$ de $B$ sur des parties
$J \subset I$. La facette fermée (ouverte) associée à $s$ est l'ensemble des
$\lambda \in C(B, \sigma)$ tels que $\lambda_{s(j)} = 1$ pour $j \in J$ et
$|\lambda_{b}| \leq 1$ (resp. $< 1$) pour $b$ au-dessus de $I \setminus J$.

\emph{Démonstration (page 24).} Soient $\lambda \neq \mu$ dans le cube. Ils sont
intérieurs à un même segment de $C$ si et seulement si la droite qui les joint
se prolonge un peu au-delà de chacun dans le cube, c'est-à-dire si, pour tout
$b$, $|t\lambda_{b} + (1-t)\mu_{b}| \leq 1$ pour $t$ un peu plus grand que 1 et
un peu plus petit que 0. Aucune coordonnée ne fait obstacle si
$\lambda_{b} = \mu_{b}$, ni si $|\lambda_{b}|, |\mu_{b}| < 1$ ; si
$\lambda_{b} \neq \mu_{b}$ et que l'un vaut $\pm 1$, la droite sort du cube de
ce côté. Donc $R(\lambda, \mu)$ équivaut à : $J(\lambda) = \{b \mid \lambda_{b}
= 1\}$ est égal à $J(\mu)$. La proposition s'ensuit : c'est le fait qu'un
produit de segments a pour facettes les produits de facettes des
segments.\footnote{Ce dernier énoncé est le nôtre. La page note
$\mathcal{R}$ et $\mathcal{B}$ par moments pour $R$ et $B$.}

\textbf{Corollaire.} L'ordre des facettes de $C(B, \sigma)$ est l'inverse de la
relation de prolongement entre sections partielles. La facette maximale
correspond à la section vide, les facettes sous-maximales aux éléments de $B$,
et les adhérences de deux facettes sous-maximales distinctes sont disjointes
--- parallèles --- si et seulement si elles correspondent à deux éléments
conjugués $b$, $-b$.\footnote{La page dit que deux facettes sous-maximales
distinctes « ne se rencontrent pas » ; deux facettes distinctes ne se
rencontrent jamais, et l'énoncé porte sur leurs adhérences. Le « ne » est un
ajout interlinéaire.}

\textbf{Exercice (page 25).} Soient $p : B \to I$ une surjection quelconque et
$\mathcal{P}$ l'ensemble des sections partielles de $p$, ordonné par
prolongement. Ses éléments maximaux forment l'ensemble
$C = \prod_{i} p^{-1}(i)$ des sections, et les successeurs de la section vide
s'identifient à $B$. Pour $b \in B$, soit $C_{b}$ l'ensemble des sections qui
prolongent $b$, c'est-à-dire telles que $s(p(b)) = b$. Pour $b \neq b'$,
\[
  C_{b} \cap C_{b'} = \emptyset \iff p(b) = p(b'),
\]
de sorte que $p(b) = p(b')$ si et seulement si $b = b'$ ou
$C_{b} \cap C_{b'} = \emptyset$. On reconstruit ainsi $B$, la relation
d'équivalence définie par $p$, donc $I$ et $p$, à isomorphisme près, à partir de
l'ensemble ordonné $\mathcal{P}$.\footnote{Une hypothèse biffée demandait des
fibres d'au moins deux éléments ; elle est en effet inutile. La surjectivité
sert à ce que $b$ et $b'$ de fibres distinctes soient prolongés par une même
section.}

\textbf{Corollaire.} Le groupoïde des surjections $B \to I$ est équivalent, par
$p \mapsto \mathcal{P}(p)$, à une sous-catégorie pleine du groupoïde des
ensembles ordonnés.\footnote{La pleine fidélité tient à ce qu'un isomorphisme
de $\mathcal{P}(p)$ est déterminé par son action sur les successeurs de la
section vide, toute section partielle étant la borne supérieure de ceux
qu'elle prolonge ; la page ne le dit pas.}

\pagerange{26}{26}
\subsection*{Le cube réel et son groupe (page 26)}

Du point de vue de la géométrie des convexes sur $\mathbf{R}$ :
\begin{enumerate}
\item[(a)] le cube standard est $C_{n} = \{x \in \mathbf{R}^{n} \mid |x_{i}| \leq 1
\ \forall i\}$ ;
\item[(b)] un cube d'un espace affine réel $E$ de dimension $n$ est une partie
$C$ telle qu'il existe un isomorphisme affine $u : \mathbf{R}^{n} \to E$ avec
$u(C_{n}) = C$ ;
\item[(c)] les $n$-cubes convexes forment un groupoïde connexe, pointé par
$(\mathbf{R}^{n}, C_{n})$, donc équivalent au groupoïde des torseurs sous
\[
  W_{n} = \operatorname{Aut}_{\mathrm{aff}}(\mathbf{R}^{n}, C_{n}).
\]
\end{enumerate}

\textbf{Proposition.} Soit $B_{n} = \{\pm e_{i} \mid 1 \leq i \leq n\}$. Alors
\[
  \operatorname{Aut}_{\mathrm{aff}}(\mathbf{R}^{n}, C_{n}) =
  \operatorname{Aut}_{\mathrm{vect}}(\mathbf{R}^{n}, B_{n}) \xrightarrow{\ \sim\ }
  \operatorname{Aut}(B_{n}, \sigma_{n}),
\]
où $\sigma_{n}$ est l'involution $b \mapsto -b$.

La démonstration est « laissée comme exercice ». Un automorphisme affine du cube
fixe son centre de symétrie, donc est linéaire ; il permute les facettes
sous-maximales $\{x_{i} = \pm 1\}$, dont les centres sont les $\pm e_{i}$ ; et
les automorphismes de $(B_{n}, \sigma_{n})$ sont les permutations signées, dont
les matrices sont orthogonales. $W_{n}$ est le groupe hyperoctaédral, d'ordre
$2^{n} n!$, groupe de Weyl des types $B_{n}$ et $C_{n}$.\footnote{La page
remarque que « $B_{n}$ est l'ens. des \emph{pts extrémaux} du convexe
$C_{n}$ » : c'est inexact, les points extrémaux de $C_{n}$ étant les $2^{n}$
points $(\pm 1, \ldots, \pm 1)$. $B_{n}$ est l'ensemble des centres des facettes
sous-maximales de $C_{n}$, c'est-à-dire, au signe de la convention près,
l'ensemble des points extrémaux du polaire $C_{n}^{\circ}$, l'hyperoctaèdre
$\{\sum |x_{i}| \leq 1\}$. La démonstration ci-dessus et les noms
« hyperoctaédral », « groupe de Weyl » sont les nôtres --- la notation $W_{n}$ est
la sienne ; pour les groupes de Weyl dans ces dossiers, voir le dossier 88.
$\sigma_{n}$ est une lecture incertaine.} En marge, la même chose en termes
d'une épure quelconque : $E = \operatorname{Hom}_{\{\pm 1\}}(B, \mathbf{R})$ et
$C = \{\lambda \mid |\lambda(b)| \leq 1 \ \forall b \in B\}$.

\textbf{Corollaires.} Le groupoïde des $n$-cubes convexes est équivalent à
celui des épures de $n$-cubes ; le groupoïde des $n$-cubes convexes orientés est
équivalent à celui des épures de $n$-cubes orientées.

Les deux groupoïdes sont connexes, et leurs groupes d'automorphismes sont
$W_{n}$ et $\operatorname{Aut}(B_{n}, \sigma_{n})$, identifiés par la
proposition. Pour les objets orientés, le sous-groupe des rotations du cube
correspond au sous-groupe des automorphismes de l'épure qui opèrent
trivialement sur $\omega_{I} \wedge \bigwedge_{i} B_{i}$ : le déterminant d'une
permutation signée est le produit de la signature de la permutation et des
signes.

\pagerange{27}{30}
\section*{III. Facettes, seconde rédaction, et polyèdres combinatoires (pages 27 à 30)}

\pagerange{27}{29}
\subsection*{Les facettes reprises (pages 27 à 29)}

La page 27 reprend la définition des facettes d'un convexe $C$ d'un espace
vectoriel réel, cette fois comme parties \emph{non vides}. Le critère qu'elle
écrit --- tout intervalle ouvert contenu dans $C$ qui rencontre $F$ est contenu
dans $F$ --- caractérise les réunions de facettes ; les facettes sont les plus
petites parties non vides qui le vérifient, et ce sont les classes de la
relation
\[
  R(x,y) : \ x = y \ \text{ou} \ x \text{ et } y \text{ sont dans un même
  intervalle ouvert contenu dans } C,
\]
la relation de la page 3. La réflexivité et la symétrie sont évidentes, la
transitivité facile, et l'on obtient que $C$ est la réunion disjointe de ses
facettes, qui sont convexes, la facette de $x$ étant la classe de
$x$.\footnote{La première rédaction de la définition, entre crochets, est
encadrée et barrée. Dans la seconde, la page écrit « si $x$ est une extrémité de
$C \cap D_{xy}$ (i.e. s'il existe $x', y'$ \ldots{} tels que $x, y$ soient
intérieurs entre $x'$ et $y'$) » : la parenthèse dit le contraire de ce qui
précède, et c'est elle qu'on suit. Telle quelle, la définition admettrait
toute réunion de facettes ; la fin de la page, « l'ens. des $y$ tels que
$R(x,y)$ est l'unique facette qui contient $x$ », montre que ce sont les classes
qui sont visées. Un croquis de triangle illustre la transitivité.}

\textbf{Proposition.} Soit $C$ un polyèdre convexe, c'est-à-dire un convexe
réunion finie d'enveloppes convexes de parties finies. Alors :
\begin{enumerate}
\item[(1)] l'ensemble des facettes est fini ;
\item[(2)] il existe une unique facette maximale, l'intérieur de $C$ dans le
sous-espace affine qu'il engendre ;
\item[(3)] pour toute facette $F$, $\overline{F}$ est une réunion de facettes, et
$F$ est l'unique facette contenue dans $\overline{F}$ de dimension maximale ;
\item[(4)] si $\dim F = d \geq 1$, $\overline{F}$ contient une facette de
dimension $d - 1$ ; et $\overline{F} = \overline{F'}$ entraîne $F = F'$ ;
\item[(5)] $C$ est l'enveloppe convexe de la réunion des facettes minimales,
c'est-à-dire de ses points extrémaux.
\end{enumerate}
Ce sont les propriétés de la page 3, pour un polytope.\footnote{Les numéros sont entourés sur la page. L'article (6), comme la fin de
(5), est écrit en oblique dans un coin et ne se lit pas au-delà de quelques
mots. En (4), la page écrit « si $\dim F = d$ » ; la facette vide étant exclue
dans cette rédaction, il faut $d \geq 1$. Une marge propose « ou encore : int. \emph{finie} de demi-espaces
fermés » comme autre définition de polyèdre ; c'est celle des secteurs
polyédraux de la page 10, et elle ne convient qu'avec la compacité : une droite
entière est une intersection finie (vide) de demi-espaces, et ne vérifie ni
(4) ni (5).}

\emph{Les facettes fermées (page 28).} Soient $C$ polyédral de dimension $d$,
$A_{C}$ le sous-espace affine qu'il engendre, $C^{0}$ son intérieur dans
$A_{C}$ --- la facette maximale ---, de sorte que $C = \overline{C^{0}}$. Les
autres facettes fermées sont les $H \cap C$, $H$ hyperplan de $A_{C}$ ne
rencontrant pas $C^{0}$ ; celles de dimension $d - 1$ correspondent
bijectivement aux $H$ tels que $H \cap C$ soit d'intérieur non vide dans $H$, et
ces intérieurs sont les facettes de dimension $d - 1$. Soit $C^{1}$ leur
réunion, et ainsi de suite. Toute facette fermée de dimension $\leq d - i$ est
de la forme $L \cap C$, $L$ sous-espace affine de $A_{C}$ de dimension $d - i$ ne
rencontrant pas $C^{0} \cup \cdots \cup C^{i-1}$ ; et parmi ces $L$, ceux pour
lesquels $L \cap C$ est d'intérieur non vide dans $L$ correspondent
bijectivement aux facettes de dimension $d - i$, qui sont ces
intérieurs.\footnote{La page affirme que les facettes fermées de dimension
$\leq d - i$ \emph{sont} les $L \cap C$ ; pour $i \geq 2$, la réciproque est
fausse. Dans un cube de dimension 3, une droite qui passe par le milieu d'une
arête, transversalement et hors du cube, ne rencontre ni l'intérieur ni les
faces carrées ouvertes, et coupe le cube en un point qui n'est pas un sommet.
La correspondance avec les facettes de dimension exactement $d - i$, elle, est
juste. On note $L$ ce que la page note $F$, réservé ici aux facettes.}

L'intersection de deux facettes fermées est une facette fermée, et :

\textbf{Proposition.} Les facettes fermées de $C$ sont les réunions de facettes
qui sont convexes et fermées.

C'est l'énoncé (6) de la page 4. La page 29 rassemble : pour une facette $F$,
\begin{enumerate}
\item[a)] $F = \operatorname{int}_{A_{F}} \overline{F}$ et $\overline{F} = A_{F}
\cap C$, où $A_{F} = A_{\overline{F}}$ ;
\item[b)] une partie $F' \subset C$ est une facette de $\overline{F}$ si et
seulement si c'est une facette de $C$ contenue dans $\overline{F}$ ;
\item[c)] $K \subset C$ est une facette fermée si et seulement si $K$ est
convexe, fermé, et réunion de facettes --- d'où la stabilité par intersection ;
\item[d)] soit $\Pi = \mathrm{Fac}(C)$, facette vide comprise. Alors 1° l'ensemble
ordonné opposé $\Pi^{\circ}$ est isomorphe à $\mathrm{Fac}(C')$ pour un polyèdre
convexe $C'$ de même dimension ; 2° pour $F' < F''$ dans $\Pi$, l'intervalle
$\Pi_{F',F''}$ est isomorphe à $\mathrm{Fac}(K)$ pour un polyèdre convexe $K$ de
dimension $\dim F'' - \dim F' - 1$.
\end{enumerate}
Ce sont les corollaires 2 et 3 des pages 11 et 12.\footnote{En b), la page écrit « facette de $F$ » pour « facette de
$\overline{F}$ ». En d), elle écrit « partie polyédrale » pour $C'$, et
« $\Pi_{F}$ » pour $\Pi_{F',F''}$. La facette vide, exclue page 27, doit être
comptée pour que 1° et 2° soient vrais : sans elle, $\Pi$ n'a pas de plus petit
élément et $\Pi^{\circ}$ pas de plus grand. Ce sont les corollaires 2 et 3 des
pages 11 et 12.}

\textbf{Corollaire.} Si $\dim F'' - \dim F' = 2$, il y a exactement deux
facettes $F$ telles que $F' < F < F''$ --- le losange du corollaire 7 de la
page 14.

\pagerange{29}{30}
\subsection*{Le polyèdre combinatoire : axiomes et recensement (pages 29 et 30)}

\textbf{Définition.} Un \emph{polyèdre combinatoire} est un ensemble ordonné fini
$\Pi$ qui vérifie :
\begin{enumerate}
\item[a)] $\Pi$ est un treillis : il a un plus petit et un plus grand élément,
et deux éléments ont une borne inférieure et une borne supérieure ;
\item[b)] $\Pi$ vérifie la condition des chaînes : deux chaînes maximales entre
deux mêmes éléments ont même longueur ;
\item[c)] si $i < j < k$, $j$ successeur de $i$ et $k$ successeur de $j$, il
existe $j' \neq j$ tel que $j'$ soit successeur de $i$ et $k$ successeur de
$j'$.
\end{enumerate}
La condition b) est la condition de Jordan--Dedekind.\footnote{La lecture de b) comme condition de Jordan--Dedekind, c'est-à-dire
comme graduation, est la nôtre ; c'est elle qu'utilise le recensement. Un petit
croquis en marge montre quatre segments issus d'un point, dans un cercle.}

Le recensement qui suit demande davantage que c) : il faut que $j'$ soit
\emph{unique}, c'est-à-dire qu'un intervalle de longueur 2 ait exactement deux
éléments intermédiaires --- le losange de la page 14, et la condition
($\Sigma$) de la page 39. Avec c) seul, un polyèdre combinatoire de dimension 1
pourrait avoir trois éléments incomparables ou plus entre $0$ et $1$. On lit
donc c) sous cette forme forte.\footnote{La page écrit « il existe $j'$ »,
sans unicité, et compte aussitôt « quatre éléments » en dimension 1.}

Il en résulte : d) tout élément $i$ de dimension $\geq 1$ est la borne
supérieure des éléments qui lui sont strictement inférieurs, et dualement tout
élément de codimension $\geq 2$ est la borne inférieure des éléments qui lui
sont strictement supérieurs. En effet, si $s = \sup\{j \mid j < i\}$ était
$< i$, $s$ serait l'unique prédécesseur de $i$, et, $s$ n'étant pas $0$, c) appliqué
à un prédécesseur de $s$ fournirait un second prédécesseur de
$i$.\footnote{La page énonce d) pour « $i$ non minimal » et « $i$ non
maximal », ajouts qui remplacent une incise biffée ; une flèche la relie à c).
Pour un sommet (successeur de $0$), les éléments strictement inférieurs se
réduisent à $0$, dont la borne supérieure n'est pas le sommet : il faut
exclure les sommets, et dualement les éléments dont $1$ est successeur.}

\emph{Recensement}, la dimension étant celle de $1$, avec $\dim 0 = -1$ :
\begin{itemize}
\item dimension $-1$ : un seul élément, $0 = 1$ ;
\item dimension $0$ : deux éléments, $0 \neq 1$ ;
\item dimension $1$ : quatre éléments, $0$, $1$ et deux éléments incomparables
$x$, $y$ --- le segment ;
\item dimension $2$ : $0$, des sommets, des arêtes, $1$ ; les sommets et les
arêtes forment un « contour » --- chaque sommet sur deux arêtes, chaque arête
sur deux sommets --- dont les composantes connexes ont au moins trois sommets ;
\item dimension $3$ : $0$, sommets, arêtes, faces, $1$. On voudrait que, pour un
sommet $x$, les éléments $y > x$, $y \neq 1$, forment un polygone, c'est-à-dire un
contour connexe, et de même, pour une face $x$, les éléments $y < x$,
$y \neq 0$ ;
\item dimension $4$ : le titre seul.
\end{itemize}

\begin{tikzcd}[column sep=small, row sep=small]
  & 1 & \\
  x \arrow[ur, no head] & & y \arrow[ul, no head] \\
  & 0 \arrow[ul, no head] \arrow[ur, no head] &
\end{tikzcd}

En dimension 2, les axiomes interdisent le digone --- deux sommets sur les deux
mêmes arêtes n'auraient pas de borne supérieure --- mais pas la réunion
disjointe de deux triangles : celle-ci est un polyèdre combinatoire qui n'est
l'ensemble des facettes d'aucun polyèdre convexe, et c'est ce qui sépare cette
notion de celle de polyèdre combinatoire strict. Les contours sont les
polygones combinatoires du dossier 89 (pages 19 et 20), non nécessairement
connexes.\footnote{À droite de la page, un hexagone dont un sommet est marqué et
un pentagone. Pour la dimension 3, la page écrit « un contour \emph{convexe} »
pour les éléments au-dessus d'un sommet, puis « un \emph{contour} » suivi de
« ($\pi_{0} = 0$) », lecture incertaine ; c'est la connexité que demandent les
conditions c) et c') de la page 32, et on l'écrit.}

\pagerange{32}{36}
\section*{IV. La dimension 2 : cartes lisses combinatoires (pages 32 à 36)}

\pagerange{32}{33}
\subsection*{Drapeaux et involutions (pages 32 et 33)}

Les pages 32 à 36 étudient un polyèdre combinatoire de dimension 3 par son
bord : les ensembles $S = \mathfrak{X}_{0}$ des sommets, $A = \mathfrak{X}_{1}$
des arêtes, $F = \mathfrak{X}_{2}$ des faces, et les incidences
$\vec{A} \subset S \times A$ (les « arcs tangents », un sommet et une arête qui
en part) et $A^{\uparrow} \subset A \times F$ (les « arcs transverses », une
arête et une face qu'elle borde).

\begin{tikzcd}[column sep=small, row sep=small]
  & \vec{A} \arrow[dl] \arrow[dr] & & A^{\uparrow} \arrow[dl] \arrow[dr] & \\
  S & & A & & F
\end{tikzcd}

Les axiomes :
\begin{enumerate}
\item[a)] $\vec{A} \to A$ est un revêtement à deux feuillets : toute arête a deux
extrémités ;
\item[a')] $A^{\uparrow} \to A$ est un revêtement à deux feuillets : toute arête
borde deux faces ;
\item[b)] l'application $\mathrm{Dr} = \vec{A} \times_{A} A^{\uparrow} \to S \times
F$ est à deux feuillets sur son image : entre un sommet et une face incidents,
il y a exactement deux arêtes ;
\item[c)] pour $s \in S$, l'ensemble $\mathfrak{X}_{\geq s}$ des arêtes et des
faces incidentes à $s$ --- un graphe, de sommets les arêtes et d'arêtes les
faces grâce à b), et un contour grâce à a') --- est un polygone, c'est-à-dire un
contour non vide et connexe ;
\item[c')] dualement, pour $f \in F$, l'ensemble $\mathfrak{X}_{\leq f}$ des
sommets et des arêtes incidents à $f$ est un polygone.
\end{enumerate}

\emph{Version commune.} Adjoignons $\mathfrak{X}_{-1} = \{0\}$ et
$\mathfrak{X}_{3} = \{1\}$. Les axiomes disent alors : pour $-1 \leq i < j \leq
3$, $j - i \geq 2$, $(i,j) \neq (-1, 3)$, $x \in \mathfrak{X}_{i}$,
$y \in \mathfrak{X}_{j}$, $x < y$, la réalisation topologique de l'intervalle
ouvert $\mathfrak{X}_{x < \cdot < y}$ est homéomorphe à la sphère
$S^{\,j - i - 2}$ --- deux points pour $j - i = 2$, un cercle pour
$j - i = 3$.\footnote{La page écrit $S^{\delta - 1}$, $\delta = j - i$ ; c'est
$S^{\delta - 2}$ : entre $0$ et une arête, $\delta = 2$, il y a deux sommets,
soit $S^{0}$. La formule de la page 38, $\delta(y) - \delta(x) - 2$, est la
bonne. « $\delta =$ » est entouré sur la page ; l'indice de $\mathfrak{X}$ est lu
d'après la définition qui le suit.} Le cas exclu $(-1, 3)$ est le bord tout
entier, qui n'est pas supposé être une sphère : les conditions disent que sa
réalisation est une surface fermée, sans point singulier --- d'où, peut-on
penser, le mot « lisse ». La page annonce que cela se généralise « aux polyèdres
lisses de dimension quelconque ».\footnote{L'interprétation de « lisse » est la
nôtre.}

\emph{Les involutions.} Par a), a') et b), un drapeau $(s, a, f)$ a un unique
voisin qui ne diffère que par le sommet, par la face, par l'arête ; d'où trois
involutions sans point fixe de $\mathrm{Dr}$, $\sigma_{1}$ (changer le sommet),
$\sigma_{2}$ (changer la face) et $\varepsilon$ (changer l'arête), avec
$\sigma_{1}\sigma_{2} = \sigma_{2}\sigma_{1} =: \sigma$, elle aussi sans point
fixe. Les conditions c) et c') assurent que $S$ et $F$ sont bien les sommets et
les faces de la carte ainsi définie sur l'ensemble $\mathrm{Dr}$.\footnote{Ce
sont les opérations $\sigma_{0}$, $\sigma_{1}$, $\sigma_{2}$ du dossier 88
(page 10) et de l'\emph{Esquisse d'un programme} : $\sigma_{1}$, $\varepsilon$,
$\sigma_{2}$ ici sont respectivement $\sigma_{0}$, $\sigma_{1}$, $\sigma_{2}$
là, et la relation $\sigma_{1}\sigma_{2} = \sigma_{2}\sigma_{1}$ est
$(\sigma_{0}\sigma_{2})^{2} = 1$. Le dictionnaire est le nôtre ; la page ne dit
pas laquelle des trois opérations change quoi, et on le lit sur les flèches
étiquetées du diagramme de la page 33.}

Soient $\rho_{s} = \varepsilon\sigma_{2}$ et $\rho_{f} = \sigma_{1}\varepsilon$,
les rotations autour d'un sommet et d'une face. On trouve (page 33) :
\begin{itemize}
\item $\widetilde{S} = \mathrm{Dr}/\langle\rho_{s}\rangle$, l'ensemble des
couples $(s, \omega)$, $s \in S$, $\omega$ l'un des deux ordres circulaires sur
les arêtes issues de $s$ qui définissent la structure de polygone de
$\mathfrak{X}_{\geq s}$ --- les \emph{sommets orientés} ;
\item $\widetilde{F} = \mathrm{Dr}/\langle\rho_{f}\rangle$, de même, les
\emph{faces orientées} ;
\item $\widetilde{A} = \mathrm{Dr}/\sigma$, l'ensemble des couples $(a, u)$,
$a \in A$, $u \in \vec{A}(a) \wedge A^{\uparrow}(a)$, où $\vec{A}(a)$ et
$A^{\uparrow}(a)$ sont les ensembles à deux éléments des arcs tangents et
transverses portés par $a$ --- les \emph{arêtes orientées}, au sens d'une
orientation locale de la surface le long de $a$.
\end{itemize}
Ce sont des revêtements doubles de $S$, $F$ et $A$.\footnote{La page note les quotients $A^{\llcorner\!\to}/\rho_{s}$ et
$A^{\llcorner\!\to}/\rho_{f}$ sans définir $\rho_{s}$, $\rho_{f}$ ; on les prend au
sens du dossier 88 et de l'\emph{Esquisse}. « Portés » est une lecture
incertaine.} Soit enfin $A^{\wedge} = \operatorname{Im}(\mathrm{Dr} \to S \times
F)$, la relation d'incidence entre sommets et faces --- les coins.

\begin{tikzcd}[column sep=large, row sep=large, nodes={font=\scriptsize}]
  & & \mathrm{Dr} \arrow[dll, "\rho_{s}"'] \arrow[dl, "\sigma_{2}"'] \arrow[d, "\sigma"] \arrow[dr, "\sigma_{1}"] \arrow[drr, "\rho_{f}"] & & \\
  \widetilde{S} \arrow[d, "\sigma_{2}"'] & \vec{A} \arrow[dl] \arrow[dr, "\sigma_{1}"] & \widetilde{A} & A^{\uparrow} \arrow[dl, "\sigma_{2}"'] \arrow[dr] & \widetilde{F} \arrow[d, "\sigma_{1}"] \\
  S & & A & & F
\end{tikzcd}

Chaque flèche étiquetée est le quotient par l'involution, ou le groupe, qui la
porte.\footnote{Sur la page, la flèche $\widetilde{S} \to S$ porte $\sigma_{1}$
et la flèche $\widetilde{F} \to F$ porte $\sigma_{2}$. Avec le sens que les
quatre autres flèches étiquetées donnent à $\sigma_{1}$ (changer le sommet) et
à $\sigma_{2}$ (changer la face), ce n'est pas possible : le passage de
$\widetilde{S}$ à $S$ oublie l'orientation en un sommet fixé, que renverse
$\sigma_{2}$ ou $\varepsilon$, non $\sigma_{1}$, qui change le sommet. On
échange les deux étiquettes. Le $\sigma_{2}$ de la flèche
$\mathrm{Dr} \to \vec{A}$ est lui-même une lecture incertaine. Sous
$\widetilde{A}$ la page écrit « $\sigma_{1} \updownarrow \sigma_{2}$ », sans
flèche vers $A$, et on n'en ajoute pas.}

\begin{tikzcd}[column sep=normal, row sep=normal]
  & \mathrm{Dr} \arrow[dl] \arrow[d] \arrow[dr] & \\
  \widetilde{S} \arrow[d] & A^{\wedge} \arrow[dl] \arrow[dr] & \widetilde{F} \arrow[d] \\
  S & & F
\end{tikzcd}

Un drapeau se projette ainsi sur un sommet orienté, un coin et une face orientée.\footnote{Le trait de $\widetilde{F}$ à $F$ n'a pas de pointe sur la page ; c'est
la projection, comme à gauche.}

On a des isomorphismes canoniques
\[
  \mathrm{Dr} \simeq \widetilde{S} \times_{S} \vec{A} \simeq A^{\uparrow}
  \times_{F} \widetilde{F} \simeq \vec{A} \times_{A} A^{\uparrow} \simeq \vec{A}
  \times_{A} \widetilde{A} \simeq \widetilde{A} \times_{A} A^{\uparrow} \simeq
  \widetilde{S} \times_{S} A^{\wedge} \simeq A^{\wedge} \times_{F} \widetilde{F}.
\]
Par exemple, un sommet orienté et une arête qui en part déterminent la face qui
suit l'arête dans l'orientation ; un coin $(s, f)$ et une orientation en $s$
déterminent celle des deux arêtes du coin qui vient la première. Un drapeau
est ainsi la même chose qu'un système d'incidences compatibles entre sommet,
arête et face orientés.

\pagerange{34}{35}
\subsection*{L'axiome D (pages 34 et 35)}

\textbf{Axiome D.} Toute partie non vide minorée de $\mathfrak{X}_{**} =
\mathfrak{X}_{0} \cup \mathfrak{X}_{1} \cup \mathfrak{X}_{2}$ a une borne
inférieure ; de façon équivalente (lemme ci-dessous), toute partie non vide
majorée a une borne supérieure. C'est dire que $\mathfrak{X}$, augmenté de $0$
et de $1$, est un treillis : l'axiome a) de la page 29.

L'ensemble étant fini, il suffit que deux éléments minorés aient une borne
inférieure. C'est automatique si l'un est un sommet, ou s'ils sont
comparables. Restent trois cas.
\begin{enumerate}
\item[1°)] \emph{Deux arêtes} ayant un sommet commun : elles n'en ont qu'un, ou
sont égales ; le graphe $(S, A)$ est simple.
\item[2°)] \emph{Deux faces} distinctes ayant un sommet commun : leurs minorants
communs ont un plus grand élément, c'est-à-dire qu'elles ont en commun soit un
seul sommet et aucune arête, soit une seule arête et ses seules extrémités.
\item[3°)] \emph{Une arête $a$ et une face $f$} non incidentes ayant un sommet
commun : la seconde extrémité de $a$ n'est pas dans $f$ ; autrement dit, une
face incidente aux deux extrémités d'une arête est incidente à l'arête.
\end{enumerate}

Regroupés, ces cas donnent les quatre conditions de la page :
\begin{enumerate}
\item[d$_{1}$)] une arête est déterminée par l'ensemble de ses extrémités ;
\item[d$_{2}$)] une face qui contient les deux extrémités d'une arête contient
celle-ci ;
\item[d$'_{1}$)] une arête est déterminée par l'ensemble de ses deux faces ;
\item[d$'_{2}$)] les deux faces d'une arête n'ont pas d'autre sommet commun que
ses extrémités ;
\end{enumerate}
d$_{1}$) et d$_{2}$) disent que toute arête est la borne supérieure de ses deux
extrémités, d$'_{1}$) et d$'_{2}$) qu'elle est la borne inférieure de ses deux
faces. À quoi il faut ajouter la partie du cas 2° que ces quatre conditions ne
couvrent pas :
\begin{enumerate}
\item[d$_{3}$)] deux faces distinctes qui ont deux sommets communs ont une arête
commune.
\end{enumerate}
L'axiome D équivaut à d$_{1}$), d$_{2}$), d$'_{1}$), d$'_{2}$) et
d$_{3}$).\footnote{La page arrête la liste à d$'_{2}$), et la page 36 appelle
« $D = $ d) + d') » leur conjonction. La condition d$_{3}$) n'en découle pas : sur
le tore $\mathbf{R}^{2}/\Lambda$, $\Lambda$ engendré par $(2,2)$ et $(3,-3)$,
le quadrillage entier donne une carte de 12 sommets, 24 arêtes et 12 faces
carrées qui vérifie a) à c') et les quatre conditions de la page, mais où les
carrés de coins $(0,0)$ et $(1,1)$ ont en commun les sommets $(1,1)$ et
$(2,2) \equiv (0,0)$, sans arête commune ; ces deux faces n'ont pas de borne
inférieure. (Vérification de l'édition, faite par un calcul.) Le cas 2° de la
page dit d'ailleurs, en partie illisible, que deux faces « qui ont en commun deux
sommets » les ont « liés par l'arête qui les joint » : c'est d$_{3}$), que le
regroupement a perdu. Le paragraphe qui précède les trois cas est encadré et
barré ; il réduisait D à « au moins trois arêtes en chaque sommet ».}

\textbf{Lemme (page 35).} Soient $\mathfrak{X}$ un ensemble ordonné, $A \subset
\mathfrak{X}$, et $B$ l'ensemble des majorants de $A$. Par définition,
$\sup A$ existe si et seulement si $B$ a un plus petit élément.
\begin{enumerate}
\item[1.] Si $A$ a un plus grand élément $\alpha$, $\sup A = \alpha$.
\item[2.] $\sup A$ existe si et seulement si $\inf B$ existe, et alors $\inf B =
\sup A$.
\end{enumerate}
La nécessité vient de 1, dualisé. Pour la suffisance, soit $\widetilde{A}$
l'ensemble des minorants de $B$ ; $A \subset \widetilde{A}$, et
$\alpha = \inf B$ est le plus grand élément de $\widetilde{A}$ ; il majore
$\widetilde{A}$, donc $A$, donc $\alpha \in B$, et il minore $B$ : c'est le plus
petit élément de $B$.\footnote{« néc. » est une lecture incertaine. Pour
$A = \emptyset$, $B = \mathfrak{X}$, et $\sup \emptyset$ existe si et seulement
si $\mathfrak{X}$ a un plus petit élément, comme la page le note.}

\textbf{Corollaire.} Pour que toute partie minorée non vide de $\mathfrak{X}$
ait une borne inférieure, il faut et il suffit que toute partie majorée non vide
ait une borne supérieure. (Si $A$ est non vide et majorée, $B$ est non vide et
minorée par les éléments de $A$.)

\pagerange{36}{36}
\subsection*{Conséquences de l'axiome D (page 36)}

\textbf{Corollaire 1.} Sous l'axiome D, tout sommet est extrémité d'au moins
trois arêtes, et toute face a au moins trois arêtes.

Soit $p(s)$ le nombre d'arêtes issues de $s$. Par c), $p(s) \neq 0$. Si
$p(s) = 1$, il n'y aurait entre $s$ et une face incidente qu'une arête, contre
b). Si $p(s) = 2$, les deux arêtes issues de $s$ seraient incidentes aux deux
mêmes faces, contre d$'_{1}$). De même pour les faces. L'inégalité $p(s) \geq 2$
ne dépend pas de D.

\textbf{Corollaire 2.} Sous l'axiome D, toute face est la borne supérieure de
l'ensemble de ses sommets, et de celui de ses arêtes ; tout sommet est la borne
inférieure de l'ensemble des faces qui y aboutissent, et de celui des arêtes
qui en partent.\footnote{Ces deux corollaires n'utilisent que d$_{1}$),
d$_{2}$), d$'_{1}$), d$'_{2}$) : une autre face contenant les sommets d'une face
$f$, au moins trois, en contiendrait les arêtes par d$_{2}$), et partagerait avec
$f$ au moins trois arêtes, contre d$'_{1}$). La démonstration est la nôtre ; la
page n'en donne pas pour le corollaire 2.}

\pagerange{37}{38}
\section*{V. Ensembles ordonnés à intervalles sphériques (pages 37 et 38)}

Soit $\mathfrak{X}$ un ensemble ordonné fini, sans hypothèse de plus petit ou de
plus grand élément, et $|\mathfrak{X}|$ sa réalisation topologique --- le
complexe simplicial des chaînes ---, avec sa décomposition canonique en les
parties fermées $|\mathfrak{X}_{\leq x}|$, qui identifie $\mathfrak{X}$ à un
ensemble ordonné de parties fermées recouvrant $|\mathfrak{X}|$. Hypothèse :

(R) pour tout $x \in \mathfrak{X}$, $|\mathfrak{X}_{< x}|$ est une sphère
topologique,

l'espace vide comptant pour la sphère de dimension $-1$. Posons
$\delta(x) = \dim |\mathfrak{X}_{<x}| + 1 = \dim |\mathfrak{X}_{\leq x}|$, la
seconde réalisation étant le cône sur la première.\footnote{(R) est la
condition qui définit aujourd'hui les ensembles ordonnés CW de Björner (1984),
ensembles des cellules des complexes CW réguliers, $|\mathfrak{X}_{\leq x}|$
étant alors une boule. Le nom est le nôtre.}

\textbf{Proposition.} $\delta$ est strictement croissante, et
$\delta(y) = \delta(x) + 1$ si et seulement si $y$ est un successeur de $x$.

Si $x < y$, $|\mathfrak{X}_{\leq x}| \subset |\mathfrak{X}_{<y}|$, d'où
$\delta(x) < \delta(y)$ ; il en résulte que $\delta(y) = \delta(x) + 1$ force $y$
à être un successeur de $x$. Inversement, si $y$ est un successeur de $x$, $x$ est
un élément maximal de $\mathfrak{X}_{<y}$ ; une chaîne maximale de
$\mathfrak{X}_{\leq x}$ est alors une chaîne maximale de $\mathfrak{X}_{<y}$, donc
un simplexe maximal de la sphère $|\mathfrak{X}_{<y}|$, de dimension
$d = \delta(y) - 1$ ; il est contenu dans $|\mathfrak{X}_{\leq x}|$, d'où
$\delta(x) \geq d$, et $\delta(x) = \delta(y) - 1$. L'argument n'utilise que la
pureté des triangulations des variétés topologiques : il suffit que
$|\mathfrak{X}_{<y}|$ soit une variété.\footnote{La page conclut de même entre
crochets, avec en marge « vérifier en général pour variété » ; la vérification
est faite : une variété topologique triangulée est pure, par l'invariance du
domaine.}

\textbf{Corollaire.} Sous (R), $\mathfrak{X}$ vérifie la condition des chaînes,
et $\delta$ est l'unique fonction $\mathfrak{X} \to \mathbf{N}$ qui augmente de 1
d'un élément à un successeur et vaut $0$ sur les éléments minimaux.

\textbf{Proposition.} Sous (R), pour $x < y$, la réalisation de l'intervalle
ouvert $\mathfrak{X}_{x < \cdot < y}$ a l'homologie d'une sphère de dimension
$\delta(y) - \delta(x) - 2$.

Elle résulte du

\textbf{Lemme.} Soit $\mathfrak{Y}$ un ensemble ordonné fini dont la
réalisation est une variété topologique de dimension $d$. Pour tous $z < x$ de
$\mathfrak{Y}$, $|\mathfrak{Y}_{>x}|$, $|\mathfrak{Y}_{<x}|$ et
$|\mathfrak{Y}_{z < \cdot < x}|$ ont l'homologie de sphères de dimensions
respectives $d - \delta(x) - 1$, $\delta(x) - 1$ et $\delta(x) - \delta(z) - 2$.

En effet, le link du sommet $x$ dans le complexe des chaînes est le joint
$|\mathfrak{Y}_{<x}| * |\mathfrak{Y}_{>x}|$, qui a l'homologie d'une sphère de
dimension $d - 1$, et un joint n'a l'homologie d'une sphère que si ses deux
facteurs l'ont ; on applique ensuite le lemme à $\mathfrak{Y} = \mathfrak{X}_{<y}$,
de dimension $\delta(y) - 1$. L'énoncé en sphères topologiques, celui de la
page, demande davantage.\footnote{La page énonce la proposition et le lemme avec « sphère
topologique », sans démonstration du lemme, et donne pour $|\mathfrak{Y}_{>x}|$
la dimension $d - \delta(x) - 2$ : c'est $d - \delta(x) - 1$ (sur le bord d'un
triangle, $d = 1$, un sommet est sous deux arêtes, soit $S^{0}$). Surtout,
l'énoncé avec des sphères topologiques est faux en général. Par le théorème de
la double suspension (Edwards, Cannon, années 1970), la double suspension
$\Sigma^{2} P$ d'une sphère d'homologie $P$ de dimension 3 non simplement
connexe est homéomorphe à $S^{5}$ ; si l'on triangule $\Sigma^{2} P$, et qu'on
prend pour $\mathfrak{X}$ l'ensemble de ses simplexes augmenté d'un plus grand
élément, (R) est vérifiée, mais l'intervalle ouvert entre un point de
suspension et le plus grand élément se réalise en $\Sigma P$, qui n'est pas une
variété. Ce qui reste vrai est l'énoncé en homologie, écrit ici, ou l'énoncé
avec des sphères pour les ensembles de facettes de polyèdres convexes, dont les
intervalles sont encore de tels ensembles (corollaire 3 de la page 12), et dans
les petites dimensions de la page 32. La démonstration du lemme est la nôtre.}

\pagerange{39}{40}
\section*{VI. Le losange et les repères (pages 39 et 40)}

Soit $\mathfrak{X}$ un ensemble ordonné fini, et $\overline{\mathfrak{X}}$
l'ensemble obtenu en lui adjoignant un plus petit élément. On suppose :

($\Sigma$) Si $x \in \overline{\mathfrak{X}}$, $y$ est un successeur de $x$ et $z$
un successeur de $y$, il existe un unique $y' \in \mathfrak{X}$, $y' \neq y$, tel
que $x < y' < z$.

C'est la propriété du losange, sous sa forme forte : l'intervalle ouvert
$]x, z[$ est exactement $\{y, y'\}$.\footnote{Deux lignes et demie barrées
suivent l'énoncé, avec un ajout interlinéaire lui-même barré. Les ensembles
ordonnés gradués dont tous les intervalles de longueur 2 sont des losanges sont
appelés aujourd'hui \emph{minces} (Björner) ; le mot est le nôtre.} L'élément
$y'$, noté $\sigma_{x,z}(y)$, est un successeur de $x$, et $z$ en est un
successeur : si l'on avait $x < y'' < y'$, on aurait $x < y'' < z$ et
$y'' \neq y'$, donc $y'' = y$, et $y < y' < z$, contre l'hypothèse que $z$ est
un successeur de $y$ ; de même de l'autre côté.

\begin{tikzcd}[column sep=small, row sep=small]
  & y' \arrow[dr] & \\
  x \arrow[ur] \arrow[dr] & & z \\
  & y \arrow[ur] &
\end{tikzcd}

\emph{Repères.} Supposons de plus que les chaînes maximales issues d'un même
élément aient toutes même longueur.\footnote{Hypothèse écrite en long dans la
marge gauche ; « issues de $x$ » est une lecture incertaine.} Un
\emph{$n$-repère} de $\mathfrak{X}$ est une suite $x_{0}, x_{1}, \ldots, x_{n}$
où $x_{0}$ est un élément minimal et $x_{i+1}$ un successeur de $x_{i}$ ; soit
$R_{n}(\mathfrak{X})$ leur ensemble --- ce sont des drapeaux. On a :
\begin{enumerate}
\item[a)] l'application d'oubli $p_{n} : R_{n} \to R_{n-1}$, pour $n \geq 1$ ;
\item[b)] des opérations $\sigma_{0}, \ldots, \sigma_{n-1}$ sur $R_{n}$, où
$\sigma_{i}$ remplace $x_{i}$ par $\sigma_{x_{i-1}, x_{i+1}}(x_{i})$, avec la
convention que $x_{-1}$ est l'élément adjoint.
\end{enumerate}
Ce sont des drapeaux munis de leurs opérations élémentaires.\footnote{La définition de $\sigma_{i}$ n'est pas écrite sur la page ; c'est la
seule que ($\Sigma$) permette, et c'est elle qui rend nécessaire l'élément
adjoint de $\overline{\mathfrak{X}}$, pour $\sigma_{0}$.}

\textbf{Propositions.}
\begin{enumerate}
\item[a)] Les $\sigma_{i}$ sont des involutions.
\item[b)] Si $0 \leq i < j \leq n - 1$ et $j \neq i + 1$, $\sigma_{i}$ et
$\sigma_{j}$ commutent.
\item[c)] Si $0 \leq i_{1} < \cdots < i_{\nu} \leq n - 1$, $\nu \geq 1$, sans deux
indices consécutifs, l'involution $\sigma_{i_{1}} \cdots \sigma_{i_{\nu}}$ est
sans point fixe.
\item[d)] Pour $p < n$, la projection $R_{n} \to R_{p}$ commute à
$\sigma_{0}, \ldots, \sigma_{p-1}$, et rend $\sigma_{p+1}, \ldots, \sigma_{n-1}$
triviales sur $R_{p}$.
\end{enumerate}
Ces énoncés sont ceux de la page.\footnote{En b), la page suppose $n \geq 3$, qui est ce qu'il faut pour que de
tels couples existent. Pour d), la fin est la ligne non barrée de la page 40,
« … en faisant que $\sigma_{n-1}, \ldots, \sigma_{p+1}$ triviales sur $R_{p}$ ».}
Les trois premiers énoncés viennent de ce que $\sigma_{i}$ ne change que la
$i$-ième place, et la change effectivement. Pour un polyèdre convexe et
$\mathfrak{X} = \mathrm{Fac}(C)$ privé de la facette vide, $\sigma_{i}$ change la
facette de dimension $i$ du drapeau : ce sont les opérations du dossier 88
(pages 6 à 8) sur les repères du complexe de Coxeter, où les produits
$\sigma_{i}\sigma_{i+1}$ ont des ordres que les pages calculent ; ici rien n'est
dit de ces produits. Pour retrouver la troisième involution $\sigma_{2}$ des
cartes de la page 32, qui change la face, il faut que $\mathfrak{X}$ contienne
son plus grand élément : les drapeaux de la carte sont alors les $3$-repères,
qui finissent tous en $1$, comme dans la « version commune » de la page 32.

La moitié supérieure de la page 40, barrée, revient aux cartes de dimension 2 :
une arête a exactement deux sommets, entre un sommet et une face il y a
exactement deux arêtes, et une condition de connexité ; en marge, des suites
$x_{0}, \ldots, x_{n}$ et quelques croquis de sommets et d'arêtes. Le dossier
s'arrête là.

\end{document}
